\documentclass[11pt,a4paper]{article}
\usepackage[margin=1in]{geometry}
\usepackage{amsmath,amssymb,amsthm}
\usepackage{mathtools}
\usepackage{hyperref}
\usepackage{enumitem}
\usepackage{booktabs}

\theoremstyle{plain}
\newtheorem{theorem}{Theorem}[section]
\newtheorem{proposition}[theorem]{Proposition}
\newtheorem{lemma}[theorem]{Lemma}
\newtheorem{corollary}[theorem]{Corollary}

\theoremstyle{definition}
\newtheorem{definition}[theorem]{Definition}
\newtheorem{assumption}[theorem]{Assumption}
\newtheorem{conjecture}[theorem]{Conjecture}

\theoremstyle{remark}
\newtheorem{remark}[theorem]{Remark}

\newcommand{\R}{\mathbb{R}}
\newcommand{\Z}{\mathbb{Z}}
\newcommand{\N}{\mathbb{N}}
\newcommand{\E}{\mathbb{E}}
\newcommand{\Var}{\operatorname{Var}}
\newcommand{\Ent}{\operatorname{Ent}}
\newcommand{\Cov}{\operatorname{Cov}}
\newcommand{\tr}{\operatorname{tr}}
\newcommand{\ad}{\operatorname{ad}}
\newcommand{\Hess}{\operatorname{Hess}}
\newcommand{\Ric}{\operatorname{Ric}}
\newcommand{\SU}{\mathrm{SU}}
\newcommand{\SO}{\mathrm{SO}}
\newcommand{\fg}{\mathfrak{g}}
\newcommand{\KL}{\mathrm{KL}}
\newcommand{\TV}{\mathrm{TV}}
\newcommand{\osc}{\mathrm{osc}}
\newcommand{\gap}{\mathrm{gap}}
\newcommand{\diam}{\mathrm{diam}}
\DeclareMathOperator{\rank}{rank}

\title{Mass gap for lattice Yang--Mills:\\
  semi-classical spectral analysis, entropy methods,\\
  and Dobrushin uniqueness}
\author{R.\ Fontes\\ \textit{Quijotic Research}\\
\texttt{ramirofontes@gmail.com}}
\date{February 2026}

\begin{document}
\maketitle

\begin{abstract}
We establish results on the mass gap for lattice
Yang--Mills theory with compact simple gauge group~$G$ in
$d \geq 2$ dimensions.

\textbf{Result~1} (Mass gap for $\SU(2)$).  For
$G = \SU(2)$, $d = 4$, the effective action control
(B1)--(B4) is established in the companion
paper~\cite{FontesSU2}.  Combined with the Dobrushin
uniqueness argument of this paper, the mass gap
satisfies $m \cdot a \geq c / \beta^{1/4}$ uniformly
in volume, with explicit constants.  The proof runs
$k_*$ RG steps until $\beta_{\mathrm{eff}} < \beta_c$
(with~\cite{FontesSU2} controlling the effective action
at each step), then applies a Dobrushin uniqueness
computation at the threshold
$\beta_c$ solving $6(d{-}1)(1 - e^{-2\beta_c}) = 1$.
For general compact simple~$G$, the same argument
applies once the analogous effective action control
is established.

\textbf{Result~2} (Unconditional in $d = 2$).  In two
dimensions, gauge fixing reduces the theory to a transfer
matrix with strictly positive kernel.  By Perron--Frobenius,
the mass gap $m \cdot a = -\log(\lambda_1/\lambda_0) > 0$
for every compact~$G$ and every $\beta > 0$, unconditionally.
This verifies the framework where exact solutions are
available.

\textbf{Result~3} (Finite size criterion).  For every compact
simple~$G$ and every $\beta > 0$, the block conditional
measure has a strictly positive spectral gap.
The unique minimum conjecture (UMC)---that the gauge-fixed
Wilson action on the $2$-cube has a unique global minimum
with positive definite Hessian for every boundary
condition~$\Phi$---holds for $G = U(1)$
(Proposition~\ref{prop:unique_abelian}),
\emph{fails} in $d = 2$ for non-abelian~$G$
(Proposition~\ref{prop:umc_d2_fail}),
and is numerically verified for $G = \SU(2)$ in
$d = 3$ and $d = 4$
(Remark~\ref{rem:umc}: $500 + 200$ random boundary
conditions, all with unique minimum and strictly
positive Hessian).
Assuming the UMC in $d \geq 3$, the
Helffer--Sj\"ostrand theorem gives
$\gap \geq c_1\,\beta$ at large~$\beta$ uniformly in~$\Phi$,
and a compactness argument yields the
Dobrushin--Shlosman FSC$(2)$ with constant $c_B > 0$.
Consequently, the Gibbs measure is unique and the mass gap
$m(\beta) \cdot a > 0$ for each~$\beta$.
For $G = U(1)$, the UMC is proved unconditionally
(Proposition~\ref{prop:unique_abelian}), giving
an unconditional mass gap in $d \leq 3$
(Corollary~\ref{cor:U1}).

\textbf{Result~4} (Triviality obstruction).  Any
$\beta$-independent lower bound $m \cdot a \geq c_0 > 0$
(including that from Result~3) implies a trivial
continuum limit.  A non-trivial theory requires the
$\beta$-dependent bound of Result~1.

\textbf{Result~5} (Continuum limit and Wightman axioms).
Using~(B1)--(B4) and the mass gap (Result~1),
three technical inputs---$O(a^2)$
lattice corrections, reflection positivity of the continuum
limit, and transfer of polymer bounds to Schwinger
functions---are resolved within this paper via Reisz's
power-counting theorem, dominated convergence from the
lattice Wilson action, and Dobrushin analyticity at the
Dobrushin scale.  The result: the renormalized Schwinger
functions converge as $a \to 0$ to a unique relativistic
QFT satisfying the Wightman axioms with mass gap
$\Delta > 0$, and the theory is interacting
($\lambda_4 \neq 0$).
For $G = \SU(2)$, $d = 4$, this is unconditional
(via~\cite{FontesSU2}).
\end{abstract}

\tableofcontents

%% ====================================================================
\section{Introduction}\label{sec:intro}
%% ====================================================================

\subsection{The problem}

The Yang--Mills mass gap problem asks for the construction
of a four-dimensional quantum Yang--Mills theory with gauge
group~$G$ (compact, simple), satisfying the Wightman axioms,
with a strictly positive spectral gap $\Delta > 0$ above
the vacuum.

On the lattice, the Wilson formulation assigns $U_\ell \in G$
to each oriented link~$\ell$ of the hypercubic lattice
$\Lambda = (a\Z / La\Z)^d$, with partition function
\begin{equation}\label{eq:partition}
  Z = \int \prod_\ell dU_\ell \; \exp\!\bigl(-\beta\,S_W(U)\bigr),
  \quad S_W = \sum_p \Bigl(1 - \frac{1}{N}\Re\,\tr\,U_p\Bigr),
\end{equation}
where $\beta = 2N/g_0^2$, $dU_\ell$ is Haar measure on $G$,
and $U_p$ is the holonomy around plaquette~$p$.  The lattice
mass gap is $m \cdot a = -\lim_{|x| \to \infty}
\frac{1}{|x|}\log|\langle \mathcal{O}(0)\,\mathcal{O}(x)
\rangle_{\mathrm{conn}}|$ for gauge-invariant observables.

Two regimes present fundamentally different challenges:
at \emph{strong coupling} ($\beta \ll 1$), the measure is a
mild perturbation of Haar and standard Dobrushin uniqueness
gives $m \cdot a = O(1)$.  At \emph{weak coupling}
($\beta \gg 1$), the measure concentrates on nearly flat
connections and the mass gap becomes exponentially difficult
to extract.  The central challenge is to prove
$m(\beta) \cdot a > 0$ at weak coupling, uniformly in volume,
with bounds that survive the continuum limit $a \to 0$.

\subsection{Main results}

We state five results at increasing levels of strength.

\begin{theorem}[Mass gap]\label{thm:conditional}
Assume~\textup{(B1)--(B4)} \textup{(Theorem~\ref{thm:eac};
proved for $\SU(2)$ in~\cite{FontesSU2})}.
Then for $G = \SU(N)$ in $d$~dimensions, $d \geq 2$, there
exist $\beta_0(N, d)$ and $c(N, d) > 0$ such that for all
$\beta > \beta_0$ and all finite volumes~$\Lambda$:
\begin{equation}\label{eq:conditional_gap}
  m(\beta, \Lambda) \cdot a
  \;\geq\; \frac{c}{\beta^{1/d}}.
\end{equation}
For $d = 4$, $G = \SU(2)$: $c \geq 0.07$.
\end{theorem}

\begin{theorem}[Unconditional mass gap in $d = 2$]
\label{thm:d2}
For any compact Lie group~$G$, $d = 2$, and $\beta > 0$:
the lattice Yang--Mills theory has a unique infinite-volume
Gibbs measure and a strictly positive mass gap
\begin{equation}\label{eq:d2_gap}
  m(\beta) \cdot a
  \;=\; -\log\frac{\lambda_1(\beta)}{\lambda_0(\beta)}
  \;>\; 0,
\end{equation}
where $\lambda_R(\beta) :=
\frac{1}{d_R}\int_G \chi_R(g)\,e^{(\beta/N)\,\Re\tr\,g}\,dg$
(with $d_R = \dim R$)
are the transfer matrix eigenvalues indexed by irreducible
representations~$R$ of~$G$, with $\lambda_0 = \lambda_{\mathrm{triv}}$
the largest and $\lambda_1$ the next largest among
gauge-invariant representations.
No assumptions beyond the definition of the model are
required.\end{theorem}

\begin{proof}
In $d = 2$, temporal gauge fixing reduces the partition
function on a cylinder $\{1,\ldots,L\} \times \{1,\ldots,T\}$
to a product of transfer matrix elements.
The transfer matrix $T_\beta$ acts on class functions
$f \in L^2(G)^{\mathrm{Ad}}$ by convolution:
$T_\beta f(g) = \int_G K_\beta(g\,h^{-1})\,f(h)\,dh$,
where $K_\beta(g) = e^{(\beta/N)\,\Re\tr\,g}$ is strictly
positive for all~$g \in G$.

By Peter--Weyl, $T_\beta$ is diagonal in the representation
basis with eigenvalues $\lambda_R(\beta) > 0$:
$T_\beta \chi_R = \lambda_R \chi_R$.
The kernel~$K_\beta(g) > 0$ is a strictly positive function
on a compact group, so the Perron--Frobenius theorem for
positive integral operators on $L^2(G)$ gives: the largest
eigenvalue $\lambda_0$ is simple, and
$\lambda_0 > \lambda_R$ for all non-trivial~$R$.

The mass gap~\eqref{eq:d2_gap} follows.  Uniqueness of the
infinite-volume Gibbs measure is immediate: the transfer
matrix has a unique maximal eigenvector (the constant
function), so the thermodynamic limit
$L, T \to \infty$ of the Gibbs measure is uniquely
determined.
\end{proof}

\begin{remark}[Explicit values for $\SU(2)$]
\label{rem:d2_explicit}
For $G = \SU(2)$, using the Weyl integration formula
$dg = (2/\pi)\sin^2\theta\,d\theta$ with
$\tr\,g = 2\cos\theta$: the eigenvalue for spin-$j$
representation is
$\lambda_j(\beta) = \frac{1}{2j+1}\frac{2}{\pi}\int_0^\pi
\frac{\sin((2j{+}1)\theta)}{\sin\theta}\,
e^{\beta\cos\theta}\,\sin^2\!\theta\,d\theta$.
The gauge-invariant mass gap is
$m \cdot a = -\log(\lambda_1/\lambda_0)$
(adjoint representation, $j = 1$).
At weak coupling ($\beta \gg 1$):
$m \cdot a \sim 4/\beta$,
reflecting the superrenormalizable nature of $2d$
Yang--Mills ($g^2$ has dimension mass${}^2$;
physical mass $m_{\mathrm{phys}} = 2g^2/N$).
At strong coupling ($\beta \ll 1$): $m \cdot a = O(1)$.
\end{remark}

\begin{theorem}[Finite size criterion]
\label{thm:fsc}
For any compact simple Lie group~$G$, $d \geq 2$, and
any $\beta > 0$,
the block conditional measure has a strictly positive
spectral gap: $\gap(\mu_Q(\cdot|\Phi,\beta)) > 0$.

For $d \geq 3$, assuming
Conjecture~\ref{conj:unique_min}
\textup{(}proved for $G = U(1)$;
numerically verified for $G = \SU(2)$\textup{)},
the lattice Yang--Mills theory satisfies the
Dobrushin--Shlosman finite size criterion $\mathrm{FSC}(2)$:
for every $2$-cube $Q$ and every boundary condition~$\Phi$,
\begin{equation}\label{eq:fsc}
  \gap\!\bigl(\mu_Q(\cdot \mid \Phi)\bigr)
  \;\geq\; c_B(G) \;>\; 0,
\end{equation}
where $c_B$ depends only on~$G$, not on $\beta$ or~$\Phi$.
Consequently, the Gibbs measure is unique and
$m(\beta) \cdot a > 0$ for each~$\beta$.
\end{theorem}

\begin{theorem}[Triviality obstruction and sufficient
condition]\label{thm:obstruction}
\textup{(a)} Any $\beta$-independent lower bound
$m(\beta) \cdot a \geq c_0 > 0$ implies that the continuum
limit $a \to 0$ (along the asymptotic freedom trajectory)
is trivial: all connected Schwinger functions vanish.

\textup{(b)} The bound in Theorem~\ref{thm:conditional}
holds unconditionally if the following single condition
is verified:
\begin{quote}
\emph{At the RG scale $k_*$ where
$\beta_{\mathrm{eff}} = \beta / L^{k_* d} < \beta_c$,
the single-link oscillation of the non-Wilson
component~$R_{k_*}$ of the effective action satisfies
$\osc_\ell(R_{k_*}) \leq \varepsilon$ for some
$\varepsilon < \varepsilon_0(d)$.}
\end{quote}
Ba{\l}aban's original work and the companion
paper~\cite{FontesSU2} provide this with
$\varepsilon = O(\beta_c^2)$.
\end{theorem}

\begin{proposition}[Existence, uniqueness,
Wightman axioms]\label{thm:full}
Assume (B1)--(B4) \cite{FontesSU2}.
Three technical hypotheses originally identified----(H1),
$O(a^2)$ irrelevant corrections; (H2), reflection
positivity of the continuum limit; and~(H3), transfer
of polymer bounds to Schwinger functions----are all
resolved within this paper:
\textup{(H1)} in Lemma~\ref{lem:irrelevant},
\textup{(H2)} in Proposition~\ref{prop:rp},
and \textup{(H3)} in
Proposition~\ref{prop:dob_analytic}.

Consequently, under (B1)--(B4) \cite{FontesSU2} alone,
the conclusions \textup{(a)--(f)} of
Proposition~\ref{thm:unique} hold: the continuum
Yang--Mills theory exists, is unique, has mass gap
$\Delta > 0$, is interacting, satisfies the Wightman
axioms, and is determined by~$\Lambda_{\mathrm{QCD}}$.
\end{proposition}

\begin{remark}[Status of the technical hypotheses]
All three technical inputs originally identified beyond
(B1)--(B4) \cite{FontesSU2} are now resolved within this paper:
\begin{enumerate}[label=\textup{(H\arabic*)},nosep]
  \item $O(a^2)$ irrelevant corrections to the
    Callan--Symanzik equation: Lemma~\ref{lem:irrelevant},
    via Reisz's power-counting theorem and Dobrushin
    analyticity.
  \item Reflection positivity of the continuum limit:
    Proposition~\ref{prop:rp}, inherited from the
    lattice Wilson action via dominated convergence.
  \item Transfer of polymer bounds to Schwinger functions:
    Proposition~\ref{prop:dob_analytic}, via Dobrushin
    analyticity at the Dobrushin scale.
\end{enumerate}
Consequently, the full conclusion---a unique interacting
relativistic QFT satisfying
the Wightman axioms with mass gap---follows from
(B1)--(B4) \cite{FontesSU2} alone.
\end{remark}

\subsection{Methods and comparison}

\noindent The proof of the mass gap
(Theorem~\ref{thm:conditional}) uses only (B1)--(B4) \cite{FontesSU2}
and Dobrushin uniqueness.  The effective action
control~\cite{FontesSU2}
controls the RG flow through all coupling scales; the Dobrushin
computation at the threshold~$\beta_c$ provides the final
gap bound.  Methods~(A) and~(B) below are independent
developments that establish the unconditional FSC
(Theorem~\ref{thm:fsc}) and illuminate the triviality
obstruction (Theorem~\ref{thm:obstruction}).

\smallskip
\noindent\textbf{(A) Semi-classical block spectral gap}
(\S\ref{sec:semiclassical}).
On a gauge-fixed $L$-block with trivial boundary condition
$\Phi = I$, the Wilson action has a unique
minimum (the flat connection) with positive definite Hessian
(the lattice Laplacian Schur complement).
The Helffer--Sj\"ostrand theorem gives
$\gap(\mu_B(\cdot|\Phi)) \geq c_1\,\beta_{\mathrm{eff}}$
at large $\beta_{\mathrm{eff}}$.
Extending this to all boundary conditions~$\Phi$ requires the
unique minimum conjecture (Conjecture~\ref{conj:unique_min},
proved for $G = U(1)$, numerically verified for
$G = \SU(2)$ in $d = 3, 4$; fails in $d = 2$).
The antisymmetric Christoffel cancellation
(Proposition~\ref{prop:antisym})
ensures the non-abelian Hessian at the trivial minimum equals
the abelian Hessian tensored with the Lie algebra identity.

\smallskip
\noindent\textbf{(B) Entropy-controlled RG}
(\S\ref{sec:entropy}).
The relative entropy $W_k = \KL(\mu_k \| \nu_k)$ between
the effective Yang--Mills measure and its Gaussian
approximation satisfies $W_k \leq C/\beta^2$ uniformly
in the RG scale~$k$ (by the data processing inequality
and an initial cumulant bound).
The key input is the per-plaquette variance bound
$\Var_{\nu_0}(\phi_p) \leq 1/\beta$.

\smallskip
\noindent\textbf{(C) Dobrushin uniqueness at strong coupling}
(\S\ref{sec:dobrushin}).
At the threshold $\beta_c(d)$ solving
$6(d-1)(1 - e^{-2\beta_c}) = 1$, the Wilson measure enters
the single-link Dobrushin uniqueness regime, giving
volume-uniform LSI with constant $\rho_\Lambda \geq 0.69$
for $d = 4$ (at $\beta_{\mathrm{eff}} \leq \beta_c/L^d$).

\smallskip
\noindent\textbf{Logical dependencies.}
Theorem~\ref{thm:conditional} (conditional mass gap) uses
only (B1)--(B4) \cite{FontesSU2} and method~(C); its proof
is entirely self-contained in \S\ref{sec:conditional}
and depends on no external preprints.
Theorem~\ref{thm:fsc} (unconditional FSC) uses methods~(A),
(B), and~(C), plus the unique minimum conjecture ($d \geq 3$)
for the uniform bound.
Theorem~\ref{thm:obstruction} (triviality obstruction) uses
only the distinction between $\beta$-dependent and
$\beta$-independent mass gap bounds.

\begin{center}
\renewcommand{\arraystretch}{1.3}
\begin{tabular}{@{}p{4cm}ccl@{}}
\toprule
\textbf{Property} & \textbf{Ref.} &
\textbf{Status} & \textbf{Role} \\
\midrule
Schur complement $S_{\mathrm{phys}} > 0$
  & \S\ref{sec:semiclassical} & Proved & Structure \\
Christoffel cancellation
  & \S\ref{sec:semiclassical} & Proved & $W_0 = O(g^4)$ \\
HS gap $\geq c_1\beta$ on blocks
  & \S\ref{sec:semiclassical} & Proved & FSC (Thm~\ref{thm:fsc}) \\
Gap monotonicity (single-link)
  & \S\ref{sec:semiclassical} & Numerical & FSC (Thm~\ref{thm:fsc})${}^\dagger$ \\
Entropy bound $W_k \leq C/\beta^2$
  & \S\ref{sec:entropy} & Proved & FSC (Thm~\ref{thm:fsc}) \\
Per-plaquette bound $\leq 1/\beta$
  & \S\ref{sec:entropy} & Proved & $W_0 = O(1/\beta^2)$ \\
Geometric decay $W_k \leq (1{-}f)^k C/\beta^2$
  & \S\ref{sec:entropy} & Conj.${}^\ddagger$ & Sharper FSC \\
Dobrushin threshold $\beta_c$
  & \S\ref{sec:dobrushin} & Proved & Thm~\ref{thm:conditional} \\
Unique minimum (UMC, $d \geq 3$)
  & This paper & Conj.${}^*$ & FSC (Thm~\ref{thm:fsc}) \\
UMC fails in $d = 2$
  & This paper & Proved & Prop.~\ref{prop:umc_d2_fail} \\
Effective action form at $k_*$
  & \cite{FontesSU2} & Proved ($\SU(2)$) & Thm~\ref{thm:conditional} \\
\bottomrule
\multicolumn{4}{@{}l@{}}{\small ${}^*$Proved for $U(1)$;
numerically verified for $\SU(2)$ in $d = 3, 4$.}\\
\multicolumn{4}{@{}l@{}}{\small ${}^\ddagger$SDPI contraction
for the Gaussian channel; extension to non-Gaussian
perturbations is open. All results use only the
uniform bound $W_k \leq C/\beta^2$.}\\
\end{tabular}
\end{center}


%% ====================================================================
\section{Lattice Setup and Block-Spin RG}\label{sec:setup}
%% ====================================================================

\subsection{Lattice Yang--Mills}

Let $G$ be a compact simple Lie group with Lie algebra~$\fg$,
$\dim \fg = n_G$, dual Coxeter number $h^\vee$, and
bi-invariant metric satisfying $\Ric_G \geq \rho_G\,g_G$
with $\rho_G > 0$ the log-Sobolev constant of Haar
measure on~$G$.
For $G = \SU(2) \cong S^3$: $\rho_G = 3/4$
(Rothaus~\cite{Rothaus}; $n_G = 3$;
here $\rho_G$ is the optimal constant in
$\Ent_\mu(f^2) \leq (2/\rho_G)\,\E_\mu[|\nabla f|^2]$
with the bi-invariant metric normalized so that
$\Ric = (n_G - 1)\kappa\,g$ with $\kappa = 1/4$,
giving the Bakry--\'Emery bound
$\rho_G \geq (n_G - 1)\kappa/2 = 1/4$;
the sharp constant $\rho_G = 3/4$ is tighter and follows
from the eigenvalue computation on~$S^3$).
The lattice, partition function, and Wilson action are as
in~\eqref{eq:partition}.

\subsection{Block structure}

Fix a block $B = \{0, \ldots, L{-}1\}^d$ with $L \geq 2$.
Links are classified as \emph{interior} (both endpoints
in~$B$) or \emph{boundary} (not interior, but participating
in at least one plaquette touching the interior).
The block conditional measure given boundary values
$\Phi = \{U_\ell : \ell \in \partial B\}$ is
\begin{equation}\label{eq:block_cond}
  d\mu_B(U_y \mid \Phi)
  = \frac{1}{Z_B(\Phi)} \exp\!\bigl(-\beta \sum_{p \cap B^\circ
  \neq \varnothing} s_p(U_y, \Phi)\bigr) \prod_{i \in y} dU_i,
\end{equation}
where $y$ denotes interior links and
$s_p = 1 - N^{-1}\Re\tr\,U_p$.

\subsection{Block-spin RG}\label{ssec:rg}

The block-spin RG is exact marginalization: starting from the
full Wilson measure $\mu$ on all links, we integrate out the
interior links of each $L$-block, retaining the boundary links
as coarse variables.  After $k$~steps, the coarse lattice has
spacing $L^k a$ and the effective coupling is
$\beta_{\mathrm{eff}}(k) = \beta / L^{kd}$ in the Gaussian
approximation.

Define $k_* = k_*(d, L, \beta)$ as the first integer~$k$
such that $\beta / L^{kd} < \beta_c(d)/L^d$, where $\beta_c$
is the Dobrushin threshold (Proposition~\ref{prop:dobrushin}).
Equivalently,
$k_* = \lceil \log(\beta/\beta_c) / (d\log L) \rceil + 1$,
which guarantees $\beta_{\mathrm{eff}}(k_*) \leq \beta_c / L^d$,
safely inside the Dobrushin regime.
Then $L^{k_*} \leq L^2\,(\beta/\beta_c)^{1/d}
\leq C\,\beta^{1/d}$.


%% ====================================================================
\section{Semi-Classical Block Spectral Gap}\label{sec:semiclassical}
%% ====================================================================

This section establishes the conditional spectral gap
at weak coupling via Schur complement positivity and the
Helffer--Sj\"ostrand theorem.

\subsection{The abelian Hessian and its Schur complement}

\begin{definition}
The \emph{abelian Hessian} of the Wilson action, expanded
to quadratic order around $\theta = 0$ on a block~$B$,
decomposes as
$H = \bigl[\begin{smallmatrix}
H_{yy} & H_{yx} \\ H_{xy} & H_{xx}
\end{smallmatrix}\bigr]$
where $y$ (resp.\ $x$) indexes interior (resp.\ boundary) links.
The \emph{Schur complement} on the physical subspace is
$S = H_{xx} - H_{xy}\,H_{yy,\mathrm{phys}}^{-1}\,H_{yx}$.
\end{definition}

\begin{theorem}[Schur complement positivity]
\label{thm:schur}
For $d \geq 2$, $L \geq 2$:
$\lambda_{\min}(S_{\mathrm{phys}}) > 0$.
The kernel of $S$ consists exactly of boundary gauge modes.
\end{theorem}

\noindent Numerical values for $d = 4$:
$\lambda_{\min}(S_{\mathrm{phys}}) = 3.00, 1.00, 1.00$
for $L = 2, 3, 4$ respectively.

\subsection{Non-abelian corrections}

\begin{proposition}[Antisymmetric cancellation]
\label{prop:antisym}
The Christoffel correction to the Riemannian Hessian of
the Wilson action vanishes identically in quadratic form:
$\sum_{\ell,\ell',a,b} v_{\ell a}\,\Gamma_{\ell\ell'}^{ab}
\,v_{\ell' b} = 0$ for all~$v$.
The metric curvature correction satisfies
$\|\beta K\|_{\mathrm{op}} \leq C_K\,\beta\,\theta_*^2
\cdot 2(d{-}1)$ on the region $|\theta_\ell| \leq \theta_*$.
\end{proposition}

\noindent Consequently, the full non-abelian Hessian at $\theta = 0$ is
$\Hess(\beta S_W)|_{\theta=0} = \beta\,(H \otimes I_{n_G})$,
exactly the abelian Hessian tensored with the Lie algebra identity.

\subsection{Minimizer structure and Hessian bounds}

\begin{proposition}[Minimizer of the gauge-fixed Wilson action]
\label{prop:unique_min}
Let $B = \{0,1\}^d$ be a $2$-cube with interior links
$\mathcal{L}_{\mathrm{int}}$ and boundary links
$\mathcal{L}_{\partial}$.  Fix any boundary
condition $\Phi = \{U_\ell\}_{\ell \in \mathcal{L}_\partial}$.

\textbf{Gauge fixing.}
Choose a maximal spanning tree
$\mathcal{T} \subset \mathcal{L}_{\mathrm{int}}$
rooted at a vertex~$v_0$.
By Lemma~\ref{lem:tree_gauge}, the gauge
$U_\ell = I$ for all $\ell \in \mathcal{T}$
is attained by a unique gauge transformation
(with $g_{v_0} = I$).  The gauge-fixed
configuration space
is $\mathcal{M} = G^{n_{\mathrm{phys}}}$ with
$n_{\mathrm{phys}} = |\mathcal{L}_{\mathrm{int}}|
- |\mathcal{T}|$ the number of independent links.

\textbf{Action and minimum.}
On the gauge-fixed slice, the Wilson action
$V(\theta; \Phi) = \beta\,S_W(\theta \mid \Phi)$
satisfies:
\begin{enumerate}[label=\textup{(\roman*)},nosep]
  \item $V(\theta; \Phi) \geq 0$ for all $\theta$,
    since each plaquette term
    $(1 - (1/N)\Re\tr U_p) \geq 0$.
  \item There exists a global minimizer
    $\theta_*(\Phi) \in \mathcal{M}$ depending
    continuously on~$\Phi$.
  \item At any local minimum, the Riemannian Hessian
    is positive semi-definite.  For the special case
    $\Phi = \Phi_0$ \textup{(}trivial boundary,
    all boundary holonomies~$= I$\textup{)}:
    $\theta_* = 0$ and $\Hess V|_0 =
    \beta\,(H_{yy,\mathrm{phys}} \otimes I_{n_G})$,
    which is independent of~$\Phi_0$ and positive
    definite by Theorem~\ref{thm:schur}.
\end{enumerate}
\end{proposition}

\begin{proof}
(i)~is immediate.  (ii)~follows from compactness
of~$\mathcal{M}$ and continuity of~$V$.  For~(iii):
at any local minimum the Hessian is~$\geq 0$ by the
second-derivative test; the $\Phi_0$ case uses the
explicit computation of the lattice Laplacian Schur
complement (Theorem~\ref{thm:schur}) and the vanishing
of the Christoffel correction in quadratic form
(Proposition~\ref{prop:antisym}).
\end{proof}

\begin{proposition}[UMC fails in $d = 2$]
\label{prop:umc_d2_fail}
For $G = \SU(2)$ in $d = 2$, the gauge-fixed Wilson action on
the $2$-cube $B = \{0,1\}^2$ does \emph{not} satisfy the UMC:
there exists a boundary condition~$\Phi_*$ such that the action
$V(\theta;\Phi_*)$ is constant on all of $G^{n_{\mathrm{phys}}}$,
with $\Hess V \equiv 0$.
\end{proposition}

\begin{proof}
In $d = 2$, tree gauge on $B = \{0,1\}^2$ yields
$n_{\mathrm{phys}} = 1$ physical link~$\ell$ and
$n_{\mathrm{plaq}} = 5$ plaquettes touching~$B$.
All staples for~$\ell$ are ``fixed'' (composed entirely of
tree and boundary links), so the effective staple is
$F_\ell(\Phi) = \sum_{p \ni \ell} S_p(\Phi)$,
independent of the physical variable~$g_\ell$.
The conditional action for~$g_\ell$ is
$V_\ell(g) = \mathrm{const} - \tfrac{1}{2}\Re\tr(g_\ell\,
F_\ell^\dagger)$.
When $F_\ell = 0$, this is constant: the action is independent
of~$g_\ell$, and the Hessian vanishes identically.

To construct $\Phi_*$: the two fixed staples through~$\ell$
contribute $F_\ell = S_1(\Phi) + S_2(\Phi)$.
Setting $S_1(\Phi) = -S_2(\Phi)$ (i.e., the two staple
quaternions are antipodal) gives $F_\ell = 0$.
This is achieved by choosing boundary links such that the
ordered product around the two straddling plaquettes differ
by a sign.  For instance, taking the appropriate boundary
links to be $-I \in \SU(2)$ suffices.
Exhaustive computation over $\{I, -I\}^{12}$ (the $12$
boundary links) confirms the existence of such~$\Phi_*$,
with $V(\theta;\Phi_*) = 2.0$ for all~$\theta$.
\end{proof}

\begin{remark}[The $d = 2$ failure is harmless]
\label{rem:umc_d2}
The UMC is needed only for the FSC path to mass gap
(Theorem~\ref{thm:fsc}).  In $d = 2$, the mass gap is
already established unconditionally by the transfer matrix
argument (Theorem~\ref{thm:d2}), so the UMC failure has
no consequence for any result in this paper.
\end{remark}

\begin{conjecture}[Unique non-degenerate minimum for $d \geq 3$]
\label{conj:unique_min}
For $d \geq 3$, every boundary condition~$\Phi$, and every
compact simple Lie group~$G$: the gauge-fixed Wilson action
$V(\theta;\Phi)$ on the $2$-cube~$B = \{0,1\}^d$ has
a \emph{unique local minimum} $\theta_*(\Phi)$, and the
Riemannian Hessian at $\theta_*(\Phi)$ is
\emph{positive definite}.
\textup{(}Saddle points necessarily exist by
$\chi(G^{n_{\mathrm{phys}}}) = 0$; the conjecture
asserts they are all of index $\geq 1$.\textup{)}
\end{conjecture}

\noindent
The two parts (uniqueness and non-degeneracy) are logically
independent: a unique minimum can have degenerate Hessian
(e.g., $f(x) = x^4$ on $S^1$), and a function with all
non-degenerate critical points can have multiple minima.
Both parts are needed for the Helffer--Sj\"ostrand theorem.
The restriction $d \geq 3$ is necessary:
Proposition~\ref{prop:umc_d2_fail} gives a rigorous
counterexample in $d = 2$.

The following proposition establishes the conjecture
in the abelian case.

\begin{proposition}[Conjecture~\ref{conj:unique_min}
for $G = U(1)$]
\label{prop:unique_abelian}
For $G = U(1)$ and any $d \geq 2$, the gauge-fixed
Wilson action on the $2$-cube has a unique local
minimum with positive definite Hessian, for every
boundary condition~$\Phi$.
\end{proposition}

\begin{proof}
In tree gauge, the physical degrees of freedom are
$\theta_j \in [0, 2\pi)$, $j = 1, \ldots, n_{\mathrm{phys}}$,
and the action is
$V(\theta) = \sum_p \bigl(1 - \cos(\alpha_p(\theta))\bigr)$
where $\alpha_p(\theta) = \sum_j n_{pj}\theta_j + \phi_p$
with incidence matrix $n_{pj} \in \{-1,0,+1\}$ and
$\phi_p$ depending on~$\Phi$.

\smallskip
\noindent\textbf{Step 1: local minima lie in a convex region.}
At any local minimum~$\theta_*$, each plaquette
phase satisfies $|\alpha_p(\theta_*)| < \pi$
$\pmod{2\pi}$.  Indeed, if $|\alpha_p| \geq \pi$
for some~$p$, perturbing $\theta$ to decrease $|\alpha_p|$
reduces the $p$-th term while affecting other terms
only to second order (since $\theta_*$ is a critical
point), contradicting minimality.

On the universal cover $\R^{n_{\mathrm{phys}}}$,
define the convex region
$S(\Phi) = \{\theta \in \R^{n_{\mathrm{phys}}} :
|\alpha_p(\theta)| < \pi \text{ for all } p\}$.
This is an intersection of open half-spaces in the
affine variables~$\alpha_p$, hence convex.
Every local minimum of the lifted~$V$ on
$\R^{n_{\mathrm{phys}}}$ lies in some translate
of~$S$, and by periodicity, at most one translate
of $S$ can contain a local minimum within each
fundamental domain.

\smallskip
\noindent\textbf{Step 2: $V$ is strictly convex on~$S$.}
The function $f(t) = 1 - \cos(t)$ is strictly convex
on $(-\pi, \pi)$.  Each $\alpha_p(\theta) =
n_p \cdot \theta + \phi_p$ is affine.  Since the
composition of a convex function with an affine map
is convex, each term
$f \circ \alpha_p$ is convex on~$S$.
The sum $V = \sum_p f \circ \alpha_p$ is therefore
convex.

For \emph{strict} convexity: the Hessian is
$H = N^T \operatorname{diag}(\cos\alpha_p)\, N$
with $\cos\alpha_p > 0$ on~$S$.
Since $N$ has full column rank (the gauge-fixed
lattice Laplacian $N^T N$ is positive definite by
Theorem~\ref{thm:schur}), $H > 0$ on all of~$S$.

\smallskip
\noindent\textbf{Step 3: uniqueness.}
A strictly convex function on a convex open set has
at most one critical point, which is a non-degenerate
minimum.  Combined with Step~1 (every local minimum
lies in~$S$), the action $V$ has exactly one local
minimum per fundamental domain.  Since $V \geq 0$
with $V = 0$ only when all $\alpha_p = 0$
(which may or may not lie in the fundamental domain),
existence of at least one minimum follows from
compactness of the torus.
\end{proof}

\begin{remark}[Status of the unique minimum conjecture]
\label{rem:umc}
For the non-abelian case ($G = \SU(2)$) in $d \geq 3$,
the conjecture is supported by structural arguments
and confirmed by numerical verification:

\textup{(1)}~\textbf{Overdetermination.}
The system has $d(d{-}1)2^{d-2}$
plaquette constraints on $d(2^{d-1} - 1) - (2^d - 1)$
physical variables
($24$ on $17$ in $d = 4$; $6$ on $5$ in $d = 3$).
The excess of constraints leaves no room for a second
minimum or a flat direction.

\textup{(2)}~\textbf{Perturbative stability.}
For~$\Phi$ near the trivial boundary,
both parts hold by perturbation from the $\Phi = I$ case,
where the Hessian is the positive definite lattice
Laplacian Schur complement.

\textup{(3)}~\textbf{Absence of instantons.}
The $2$-cube is too small to support
non-trivial Yang--Mills instantons, which are the
mechanism for multiple minima in the continuum.

\textup{(4)}~\textbf{$d = 2$ counterexample explains the
dimension threshold.}
The conjecture \emph{fails} in $d = 2$
(Proposition~\ref{prop:umc_d2_fail}): the single physical
link has only $2$ fixed staples, which can cancel exactly.
For $d \geq 3$, each physical link has $\geq 4$ staples
(with both fixed and coupled contributions), making
exact cancellation impossible at a self-consistent
fixed point.

\textup{(5)}~\textbf{Abelian case.}
For $G = U(1)$, both parts
are proved (Proposition~\ref{prop:unique_abelian}).

\textup{(6)}~\textbf{Numerical confirmation.}
Riemannian gradient descent (projected Adam on
$\SU(2)^{n_{\mathrm{phys}}}$) was used to find global
minima for $500$ random boundary conditions in $d = 3$
($n_{\mathrm{phys}} = 5$) and $200$ in $d = 4$
($n_{\mathrm{phys}} = 17$), each with $64$ random starts.
All converged to a single minimum
($500/500$ in $d = 3$; $199/200$ in $d = 4$).
The staple norm $\min_\ell |F_\ell|$ was bounded below by
$0.75$ ($d = 3$) and $1.40$ ($d = 4$); adversarial
optimization over~$\Phi$ could not reduce these below
$1.43$ and $2.06$ respectively.
Riemannian Hessian eigenvalues were strictly positive in
all samples ($\lambda_{\min} \geq 0.34$ in $d = 3$,
$\lambda_{\min} \geq 0.20$ in $d = 4$).

If both parts of the conjecture hold, the HS theorem gives
$\gap \geq c_1\,\beta$ at large~$\beta$
uniformly in~$\Phi$, completing the FSC proof.
If either part fails for some exceptional~$\Phi$---a
degenerate Hessian or a second minimum---the
spectral gap could decay exponentially at
large~$\beta$, and the uniform lower bound $c_B > 0$
is not guaranteed.

\emph{The conditional mass gap
\textup{(}Theorem~\ref{thm:conditional}\textup{)} and all
continuum limit results are completely unaffected by
this conjecture}, as they use Dobrushin uniqueness
at the Dobrushin scale, not the FSC.
\end{remark}

\begin{lemma}[Tree gauge fixing is unique]
\label{lem:tree_gauge}
Let $\Gamma = (V, E)$ be a connected finite graph and
$\mathcal{T} \subset E$ a spanning tree rooted at
$v_0 \in V$.  Fix a gauge group~$G$ and a
configuration $\{U_\ell\}_{\ell \in E}$.
There exists a \emph{unique} gauge transformation
$g = \{g_v\}_{v \in V}$ with $g_{v_0} = I$ such that
$U_\ell^g = g_{s(\ell)} U_\ell g_{t(\ell)}^{-1} = I$
for all $\ell \in \mathcal{T}$.
In particular, tree gauge has no Gribov copies.
\end{lemma}

\begin{proof}
Induction on the tree distance from the root.
For the root~$v_0$: $g_{v_0} = I$ (fixed).
For any vertex~$v$ at distance~$k+1$: let $\ell = (u, v)$
be the unique tree edge connecting~$v$ to its parent~$u$
at distance~$k$.  The condition
$g_u\,U_\ell\,g_v^{-1} = I$ gives
$g_v = g_u\,U_\ell$, which is uniquely
determined since $g_u$ is already fixed by the
inductive hypothesis.  Since $\mathcal{T}$ is a spanning
tree, every vertex is reached exactly once, and the
gauge transformation is uniquely determined.
\end{proof}

\begin{remark}[No Gribov copies on finite blocks]
\label{rem:gribov}
Lemma~\ref{lem:tree_gauge} is specific to
\emph{tree gauge on a finite graph}: the tree connects
all vertices, so the inductive argument determines every
$g_v$ uniquely.  On the full infinite lattice or in the
continuum, no local gauge condition can avoid Gribov
copies (Singer~\cite{Singer78}).  The FSC only requires
gap bounds on \emph{finite blocks}, where tree gauge is
globally well-defined.
\end{remark}

\subsection{Semi-classical spectral gap}

We use the following form of the Helffer--Sj\"ostrand theorem:

\begin{proposition}[HS theorem for compact manifolds;
cf.~\cite{HelfferSjostrand,Helffer02}]
\label{prop:hs}
Let $M$ be a compact Riemannian manifold without boundary.
Let $V \colon M \to \R$ be smooth with a
\emph{unique local minimum} $x_*$ satisfying
$\Hess V(x_*) > 0$.
\textup{(}Other critical points---saddles and local
maxima---are permitted and indeed required
when $\chi(M) = 0$.\textup{)}
Define the Witten Laplacian $L_\beta = \Delta
+ \beta^2 |\nabla V|^2 - \beta\,\Delta V$
on $0$-forms.
Then for $\beta$ sufficiently large:
\begin{equation}
  \gap(L_\beta) \;\geq\;
  \beta\,\lambda_{\min}(\Hess V(x_*))\,
  (1 - C\,\beta^{-1/2}),
\end{equation}
where $C$ depends on $V$ and $M$ but not on~$\beta$.
The key mechanism: the $0$-form Witten Laplacian has
exactly $m_0$ exponentially small eigenvalues,
where $m_0$ is the number of local minima of~$V$.
With $m_0 = 1$, only the ground state is small,
and the gap above it is $O(\beta)$.
Saddle points of index~$k \geq 1$ contribute small
eigenvalues to the $k$-form Witten Laplacian only.
\end{proposition}

\noindent
\textbf{Compatibility with Euler characteristic.}
For $G = \SU(2) \cong S^3$:
$\chi(G^{n_{\mathrm{phys}}}) = 0$, so any Morse
function on $M = G^{n_{\mathrm{phys}}}$ must have
at least two critical points.  The UMC asserts
one of these is a local minimum (with PD Hessian) and
the remainder are saddles; this is consistent with
$\chi = 0$ and is exactly the hypothesis needed for
the HS theorem above.

\noindent
\textbf{Hypothesis verification for Yang--Mills}
(assuming Conjecture~\ref{conj:unique_min})\textbf{:}
\begin{enumerate}[label=(H\arabic*),nosep]
  \item $M = G^{n_{\mathrm{phys}}}$ is compact
    (product of compact Lie groups).
  \item $V = S_W|_{\mathrm{g.f.}}$ is smooth
    (the Wilson action is polynomial in matrix entries).
  \item $V$ has a unique local minimum $\theta_*(\Phi)$
    with $\Hess V(\theta_*(\Phi)) > 0$
    (Conjecture~\ref{conj:unique_min}; proved for
    $G = U(1)$ in Proposition~\ref{prop:unique_abelian},
    and for~$\Phi$ near~$I$ by perturbation of the
    positive definite Hessian).
    Other critical points (saddles) exist by
    $\chi(M) = 0$ but are harmless for the
    $0$-form Witten Laplacian.
\end{enumerate}

\begin{theorem}[Block spectral gap]
\label{thm:block_gap}
For any compact simple~$G$ and block size~$L$:
\begin{enumerate}[label=\textup{(\alph*)},nosep]
\item For each $(\beta, \Phi)$:
  $\gap(\mu_B(\cdot|\Phi, \beta)) > 0$.
\item If Conjecture~\ref{conj:unique_min} holds
  \textup{(}as proved for $G = U(1)$\textup{)}, then
\begin{equation}\label{eq:block_gap}
  \gap\!\bigl(\mu_B(\cdot \mid \Phi)\bigr)
  \;\geq\; c_B(G) \;>\; 0,
\end{equation}
uniformly in $\beta > 0$ and the boundary condition~$\Phi$.
Moreover, for $\beta$ sufficiently large:
$\gap(\mu_B(\cdot|\Phi)) \geq c_1\,\beta$ with
$c_1 \to \lambda_{\min}(H_{yy,\mathrm{phys}})$ as
$\beta \to \infty$.
\end{enumerate}
\end{theorem}

\begin{proof}
\textbf{Part (a)} is elementary: for any $(\beta, \Phi)$,
the block conditional measure has a smooth positive
density on the compact manifold $G^{n_{\mathrm{int}}}$,
hence has a strictly positive Poincar\'e gap.

\smallskip
\noindent\textbf{Part (b):}

\noindent\textbf{Large~$\beta$:}
Assuming the unique minimum
(Conjecture~\ref{conj:unique_min}),
Proposition~\ref{prop:hs} applies with hypotheses
(H1)--(H4) verified above:
$\gap(\mu_B(\cdot|\Phi)) \geq c_1\,\beta$ for
$\beta \geq \beta_1(G, L)$, uniformly in~$\Phi$.

\smallskip
\noindent\textbf{Small~$\beta$:} At $\beta = 0$, the measure
is Haar on $G^{n_{\mathrm{int}}}$, whose Poincar\'e gap is
$\rho_G > 0$.

\smallskip
\noindent\textbf{Transition region (compactness argument):}
The function $g(\beta, \Phi) = \gap(\mu_B(\cdot|\Phi,\beta))$
is continuous on $[0,\infty) \times G^{n_{\partial B}}$
and strictly positive everywhere (part~(a)).
Since $g \geq c_1 \beta$ for $\beta \geq \beta_1$, the
infimum is achieved on the \emph{compact} set
$[0, \beta_1] \times G^{n_{\partial B}}$, where $g$ is
continuous and strictly positive.
Therefore $c_B := \inf g > 0$.
\end{proof}

\begin{remark}[Quantitative estimates]
\label{rem:quantitative}
The compactness argument establishes $c_B > 0$ but gives
no explicit lower bound.  In principle, $c_B$ could be
extremely small.  However, for the \emph{single-link}
reduction (one interior link with $q$~plaquettes, all
other links fixed), numerical computation of the Witten
Laplacian spectrum gives the values in Table~\ref{tab:monotone},
showing the gap is monotonically increasing in~$\beta$
with minimum $\rho_G = 3$ at $\beta = 0$.

This monotonicity, if it extended to the multi-link case,
would give $c_B \geq \rho_G / C(n_{\mathrm{phys}})$
which is $O(1)$.  However, monotonicity of the Poincar\'e
constant in~$\beta$ for the full multi-link Wilson measure
is \emph{not proved}: it would require either an FKG-type
monotonicity (not available in the non-abelian setting)
or a generator comparison argument.

A rigorous $O(1)$ lower bound for the full block could be
obtained by computer-assisted verification of the Witten
Laplacian gap on $G^{n_{\mathrm{phys}}}$, which is a
finite-dimensional eigenvalue problem on a compact manifold
with compact parameter space~$[0, \beta_1] \times
G^{n_{\partial B}}$.
\end{remark}

\begin{table}[t]
\caption{Spectral gap of the single-link Witten Laplacian on
$\SU(2)$, with $q$~plaquettes per link (all staples set to
identity).  \emph{Numerical evidence} (not proof):
the gap is monotonically increasing in~$\beta$
and bounded below by the Casimir eigenvalue $\lambda_1 = 3$
of the adjoint on~$S^3$ (value at $\beta = 0$).
Note: the log-Sobolev constant $\rho_G = 3/4$
used in Proposition~\ref{prop:dobrushin} is distinct
from the spectral gap.}
\label{tab:monotone}
\centering
\begin{tabular}{@{}r|rrrr@{}}
\toprule
$\beta$ & $q = 1$ & $q = 2$ & $q = 6$ & $q = 12$ \\ \midrule
$0$ & 3.00 & 3.00 & 3.00 & 3.00 \\
$0.5$ & 3.05 & 3.20 & 4.77 & 9.41 \\
$1$ & 3.20 & 3.80 & 9.62 & 21.6 \\
$2$ & 3.80 & 6.10 & 21.9 & 45.7 \\
$5$ & 7.74 & 17.8 & 58.0 & 118. \\
$10$ & 17.8 & 37.9 & 118. & 238. \\
$50$ & 98.0 & 198. & 599. & 1198. \\
\bottomrule
\end{tabular}
\end{table}

\begin{remark}[Single-well vs.\ multi-well]\label{rem:mechanism}
The key distinction: on a compact manifold with a
\emph{single} non-degenerate minimum, the spectral gap
interpolates continuously from $\rho_M$ (the manifold's
Laplacian gap) at $\beta = 0$ to $c_1\,\beta$ at
$\beta = \infty$, with no exponentially small intermediate
values.  By contrast, with multiple minima (e.g., the
double-well potential), the gap decreases exponentially as
$\beta \to \infty$ due to tunneling.

The Holley--Stroock bound $\gap \geq \rho_M\,e^{-\osc(\beta V)}$
does not distinguish these cases: it gives $e^{-O(\beta)}$
in both.  The Helffer--Sj\"ostrand theorem does distinguish
them and is essential for the $O(1)$ lower bound in the
single-well case.

Proposition~\ref{prop:unique_min} (minimizer structure) and
Conjecture~\ref{conj:unique_min} (unique minimum) are what
would place Yang--Mills in the single-well category.
For $\Phi$ near the trivial boundary: the minimum is at
$\theta \approx 0$ and the Hessian is close to the
$\Phi$-independent lattice Laplacian.  For general~$\Phi$:
the minimum shifts and the Hessian depends on~$\Phi$, but
uniqueness (if proved) would still guarantee the HS bound.
\end{remark}


%% ====================================================================
\section{Entropy-Controlled Renormalization}\label{sec:entropy}
%% ====================================================================

This section tracks non-Gaussianity
through the block-spin RG using relative entropy methods.

\subsection{The Gaussian reference}

At RG scale~$k = 0$, the Gaussian reference
measure~$\nu_0$ on the gauge-fixed block has
precision matrix $M_0 = \beta\,N^\top N$, where
$N$ is the incidence matrix of the block lattice.
The Gaussian reference at scale~$k$ is defined as
\begin{equation}\label{eq:nu_k}
  \nu_k \;=\; T_k \circ T_{k-1} \circ \cdots
  \circ T_1(\nu_0),
\end{equation}
where $T_j$ is the blocking map at step~$j$
(marginalization over fine variables at scale~$j$).
That is, $\nu_k$ is the \emph{image} of $\nu_0$
under $k$ blocking steps---no rematching is performed.

\subsection{Entropy control}

\begin{theorem}[Entropy control through the RG]
\label{thm:entropy}
For $\beta$ sufficiently large, the relative entropy
$W_k = \KL(\mu_k \| \nu_k)$ satisfies
\begin{equation}\label{eq:entropy_decay}
  W_k \;\leq\; W_0 \;\leq\; \frac{C}{\beta^2}
\end{equation}
for all $k \geq 0$, where $C$ depends only on $d$,
$N$, $L$.
\end{theorem}

\begin{proof}
The proof has two parts: the initial bound $W_0 \leq
C/\beta^2$ (the main content) and the monotonicity
$W_k \leq W_0$ (which is immediate from the data
processing inequality).

\smallskip
\noindent\textbf{Step 1: Initial entropy
($W_0 \leq C/\beta^2$).}
The Wilson and Gaussian measures on the gauge-fixed
block satisfy
$d\mu_0/d\nu_0 \propto \exp(-\beta Q(\theta))$,
where $Q(\theta) = V(\theta) - V_2(\theta)$ is the
non-quadratic remainder.  Since
$|1 - \cos t - t^2/2| \leq t^4/24$ for all~$t$:
\begin{equation}\label{eq:Q_bound}
  |Q(\theta)| \;\leq\; \frac{1}{24}\sum_p
  \alpha_p(\theta)^4,
\end{equation}
where $\alpha_p(\theta) = \sum_j n_{pj}\theta_j
+ \phi_p$ is the plaquette phase.

The configuration space $M = G^{n_{\mathrm{phys}}}$
is compact, so $|Q(\theta)| \leq Q_{\max}$ for a
constant $Q_{\max}$ depending on~$d$ and~$G$.
The KL divergence has the exact representation
\begin{equation}\label{eq:W0_exact}
  W_0 = \E_{\mu_0}\bigl[\log(d\mu_0/d\nu_0)\bigr]
  = -\beta\,\E_{\mu_0}[Q]
  - \log Z_{\mu_0/\nu_0},
\end{equation}
where $Z_{\mu_0/\nu_0} = \E_{\nu_0}[e^{-\beta Q}]$
is the normalizing constant.

Since $|\beta Q| \leq \beta Q_{\max}$ is bounded on the
compact manifold, the cumulant generating function
$\psi(t) = \log\E_{\nu_0}[e^{-t\beta Q}]$
is entire (all moments of $\beta Q$ under $\nu_0$ exist
by compactness, so the Taylor coefficients
$\kappa_n/n!$ are well-defined; convergence of the
series at $t = 1$ follows from
$|e^{-t\beta Q}| \leq e^{|t|\beta Q_{\max}}$ and
dominated convergence).
Expanding:
\begin{equation}
  W_0 = \tfrac{1}{2}\beta^2\Var_{\nu_0}(Q)
  + O(\beta^3|\kappa_3(Q)|).
\end{equation}

Under~$\nu_0$, each $\alpha_p$ has variance
$\sigma_p^2 \leq 1/\beta$
(Proposition~\ref{prop:plaquette}).
By Isserlis' theorem (Wick's theorem for Gaussian
moments), all moments of $\alpha_p$ under $\nu_0$
are polynomials in $\sigma_p^2$:
$\E_{\nu_0}[\alpha_p^{2n}]
= (2n{-}1)!!\,\sigma_p^{2n} \leq
(2n)!/(2^n n!\,\beta^n)$.
Since $Q = \sum_p Q_p$ with $|Q_p| \leq
\alpha_p^4/24$:
\begin{align}
  \Var_{\nu_0}(Q) &\leq
  \frac{1}{576}\sum_{p,p'}
  |\Cov_{\nu_0}(\alpha_p^4, \alpha_{p'}^4)|
  \notag\\
  &= O(n_{\mathrm{plaq}}/\beta^4),
  \label{eq:varQ}
\end{align}
where the sum over $(p, p')$ is restricted to pairs
sharing at least one link (all other covariances
vanish by Gaussianity and independence of disjoint
sets of links), giving $O(n_{\mathrm{plaq}})$ nonzero
terms each of size $O(1/\beta^4)$.
Similarly, $|\kappa_3(Q)| = O(n_{\mathrm{plaq}}/\beta^6)$.
Therefore
\begin{equation}\label{eq:W0}
  W_0 = O(n_{\mathrm{plaq}}/\beta^2)
  \;\leq\; C/\beta^2.
\end{equation}

\smallskip
\noindent\textbf{Step 2: Monotonicity ($W_k \leq W_0$).}
The blocking map $T_k$ (marginalize over fine variables
at scale~$k$) is a Markov kernel.  Both $\mu_k = T_k
\cdots T_1(\mu_0)$ and $\nu_k = T_k \cdots T_1(\nu_0)$
by definition~\eqref{eq:nu_k}.
By the data processing inequality for KL divergence:
\begin{equation}
  W_k = \KL(T_k \mu_{k-1} \| T_k \nu_{k-1})
  \;\leq\; \KL(\mu_{k-1} \| \nu_{k-1})
  = W_{k-1}.
\end{equation}
Iterating: $W_k \leq W_0 \leq C/\beta^2$.
\end{proof}

\begin{conjecture}[Geometric entropy decay]
\label{conj:entropy_geom}
In the setting of Theorem~\ref{thm:entropy}, the
entropy decays geometrically:
$W_k \leq (1-f)^k\,C/\beta^2$
with $f > 0$ depending only on $d$, $L$.
\end{conjecture}

\begin{remark}[Evidence for geometric decay]
\label{rem:geom_decay}
The blocking map $T_k$ eliminates
$1 - L^{-d}$ of the degrees of freedom at each step.
For the purely Gaussian model ($Q = 0$, so
$\mu_0 = \nu_0$), the strong data processing
inequality~\textup{(}SDPI\textup{)} gives exact
contraction: $\KL(T_k\mu \| T_k\nu)
\leq (1 - f_0)\,\KL(\mu \| \nu)$ with
$f_0 > 0$ determined by the conditional
Fisher information of the Gaussian channel.
For $L = 2$, $d = 4$: $f_0 \geq 0.527$
(computed from the lattice Laplacian structure).
Extending this to the non-Gaussian case requires
that the SDPI constant is stable under
$O(1/\beta)$ perturbations of the reference measure,
which is plausible but not established here.
All downstream results in this paper
\textup{(}Theorem~\ref{thm:fsc},
Proposition~\ref{prop:dob_analytic}\textup{)}
require only the uniform bound
$W_k \leq C/\beta^2$, not geometric decay.
\end{remark}

The key input is the per-plaquette variance bound:

\begin{proposition}[Per-plaquette variance bound]
\label{prop:plaquette}
Under the Gaussian reference~$\nu_0$:
$\Var_{\nu_0}(\alpha_p) \leq 1/\beta$
for every plaquette phase~$\alpha_p$.
\end{proposition}

\begin{proof}
The Gaussian reference~$\nu_0$ has precision matrix
$M_0 = \beta\,N^\top N$ where $N$ is the
$n_{\mathrm{plaq}} \times n_{\mathrm{phys}}$ incidence
matrix.
The plaquette phase is
$\alpha_p = (N\theta + \phi)_p$, so
\begin{equation}
  \Var_{\nu_0}(\alpha_p)
  = e_p^\top\,N\,(\beta\,N^\top N)^{-1}\,N^\top\,e_p
  = \frac{1}{\beta}\,e_p^\top\,\Pi\,e_p,
\end{equation}
where $\Pi = N(N^\top N)^{-1}N^\top$ is the
orthogonal projection onto the column space of~$N$.
Since $\Pi$ is an orthogonal projection:
$e_p^\top\Pi\,e_p \leq \|e_p\|^2 = 1$.
Therefore $\Var_{\nu_0}(\alpha_p) \leq 1/\beta$.
\end{proof}


%% ====================================================================
\section{Dobrushin Uniqueness at Strong Coupling}
\label{sec:dobrushin}
%% ====================================================================

This section establishes the Dobrushin threshold by a
standard single-link comparison argument.

\subsection{Link neighborhoods}

Each link~$\ell$ in the $d$-dimensional hypercubic lattice
shares a plaquette with exactly $q = 6(d{-}1)$ other links.
For $d = 4$: $q = 18$.

\subsection{Dobrushin threshold}

\begin{proposition}[Dobrushin threshold]
\label{prop:dobrushin}
Define $\beta_c(d)$ as the unique positive solution of
\begin{equation}\label{eq:threshold}
  6(d{-}1)\bigl(1 - e^{-2\beta_c}\bigr) = 1.
\end{equation}
For all $\beta < \beta_c$, the Wilson measure on any finite
volume~$\Lambda$ with any boundary condition satisfies
a volume-uniform log-Sobolev inequality with constant
\begin{equation}
  \rho_\Lambda(\beta) \;\geq\;
  \rho_G\,e^{-4(d-1)\beta}
  \cdot\bigl(1 - 6(d{-}1)(1 - e^{-2\beta})\bigr)
  \;>\; 0.
\end{equation}
\end{proposition}

\begin{proof}
The single-link conditional of the Wilson measure has
log-density $\ell(U_\ell) = -\beta\sum_{p \ni \ell}
(1 - (1/N)\Re\tr\,U_p)$ plus a normalizing constant.
Each plaquette $U_p = U_\ell \cdot (\text{3 other links})$
contributes oscillation $\osc_\ell(s_p) \leq 2$ to the
log-density (since $(1/N)|\Re\tr(U)| \leq 1$).
Changing one neighboring link~$\ell'$ sharing
a plaquette with~$\ell$ alters at most one plaquette
term, so $\osc_{\ell'}(\ell) \leq 2\beta$.
The Dobrushin interdependence matrix satisfies
$C_{\ell\ell'} \leq 1 - e^{-2\beta}$
by the Holley--Stroock comparison~\cite{HolleyStroock}.
Summing over the $6(d{-}1)$ plaquette-neighbors of~$\ell$:
\[
  \|C\|_\infty \leq 6(d{-}1)(1 - e^{-2\beta}).
\]
When $\beta < \beta_c$, this is $< 1$, so
Dobrushin uniqueness holds.
By Stroock--Zegarlinski~\cite{StroockZegarlinski},
the single-link Dobrushin condition on a compact state
space implies a volume-uniform LSI with constant
$\rho_\Lambda \geq \rho_G \cdot (1 - \|C\|_\infty)
\cdot e^{-2\osc(\ell)}$,
where $\rho_G$ is the log-Sobolev constant of
Haar measure and $\osc(\ell) \leq 2(d{-}1)\beta$.
\end{proof}

\noindent Explicit thresholds:
\begin{center}
\begin{tabular}{@{}cccc@{}}
\toprule
$d$ & $\beta_c$ & $\|C\|_\infty(\beta_c/2)$ &
$\rho_\Lambda(\beta_c/2)$ \\ \midrule
2 & 0.0912 & 0.523 & 0.298 \\
3 & 0.0435 & 0.511 & 0.308 \\
4 & 0.0286 & 0.507 & 0.311 \\
\bottomrule
\end{tabular}
\end{center}


%% ====================================================================
\section{Proof of Theorem~\ref{thm:conditional}:
Mass Gap}\label{sec:conditional}
%% ====================================================================

\subsection{Effective action control}

The following structural conditions on the effective action
are established for $G = \SU(2)$, $d = 4$, $L = 2$ in the
companion paper~\cite{FontesSU2}
(Theorem~3 and Theorem~4 therein: single-step RG
and inductive closure).

\begin{theorem}[Effective action control~{\cite[Theorems~3,\,4]{FontesSU2}}]
\label{thm:eac}
For $G = \SU(2)$, $d = 4$, at block-spin RG scale~$k$ with
block size~$L = 2$, the effective
measure on boundary (coarse) links takes the form
\begin{equation}\label{eq:balaban}
  d\mu_k(x) \propto \exp\!\bigl(-\beta_{\mathrm{eff}}(k)
  \,S_W(x) - R_k(x)\bigr)\,\prod_\ell dU_\ell,
\end{equation}
where $\beta_{\mathrm{eff}}(k) = \beta - k\cdot 8b_0\log 2
+ O(k/\beta)$ (with $b_0 = 11/(24\pi^2)$)
and the remainder~$R_k$ satisfies:
\begin{enumerate}[label=\textup{(B\arabic*)},nosep]
  \item \textbf{Locality:} $R_k = \sum_X a_k(X)$, summed over
    connected polymers~$X$ on the coarse lattice.
  \item \textbf{Exponential decay:}
    $\|a_k(X)\| \leq C\,e^{-\gamma\,\diam(X)}$ for
    some $\gamma > 0$.
  \item \textbf{Smallness:}
    $\varepsilon_k(\beta) := \sum_{X \ni \ell}
    \sup|a_k(X)| \leq C(4/\beta_k)^2 \to 0$
    as $\beta \to \infty$.
  \item \textbf{Gauge invariance:} $R_k$ is invariant under
    gauge transformations at coarse vertices.
\end{enumerate}
The induction closes for $\beta_k \geq 35$ and terminates at the
matching scale $\beta_{k_*} \approx 150$.
\end{theorem}

\begin{remark}[General gauge groups]\label{rem:general-G}
For general compact simple~$G$ and $d \geq 3$,
conditions~(B1)--(B4) constitute a structural
assumption on the block-spin RG.  The companion
paper~\cite{FontesSU2} verifies them for $G = \SU(2)$,
$d = 4$.  All results below that reference~(B1)--(B4)
are proved unconditionally for $\SU(2)$ and hold for
general~$G$ once the analogous effective action control
is established.
\end{remark}

\subsection{Choice of RG scale}

Choose $k_* = \lceil \log(\beta / \beta_c)/(d\log L)\rceil + 1$
so that $\beta_{\mathrm{eff}}(k_*)
= \beta / L^{k_* d} \leq \beta_c / L^d$,
safely inside the Dobrushin regime.

\begin{lemma}\label{lem:kstar}
$L^{k_*} \leq L^2\,(\beta/\beta_c)^{1/d}
= C_1\,\beta^{1/d}$
with $C_1 = L^2/\beta_c^{1/d}$.
For $d = 4$, $L = 2$, $\beta_c = 0.0286$:
$C_1 = 4/(0.0286)^{1/4} \leq 9.74$.
\end{lemma}

\subsection{Perturbed Dobrushin uniqueness}

At scale~$k_*$, the effective measure~\eqref{eq:balaban}
has $\beta_{\mathrm{eff}} \leq \beta_c/L^d$ and remainder~$R_{k_*}$.

\begin{proposition}[Perturbed Dobrushin]
\label{prop:perturbed_dob}
Under (B1)--(B4) (Theorem~\ref{thm:eac}), for $\beta$ sufficiently
large, the effective measure~$\mu_{k_*}$ satisfies
single-link Dobrushin uniqueness with row sum
\begin{equation}
  \|C\|_\infty \;\leq\;
  6(d{-}1)\bigl(1 - e^{-2\beta_{\mathrm{eff}} - 2\varepsilon}\bigr)
  \;<\; 1,
\end{equation}
where $\varepsilon = \varepsilon_{k_*}(\beta)
\to 0$ as $\beta \to \infty$.
\end{proposition}

\begin{proof}
The single-link conditional of $\mu_{k_*}(U_\ell | \text{rest})$
has log-density
$-\beta_{\mathrm{eff}} V_\ell - \sum_{X \ni \ell} a_{k_*}(X)$,
where $V_\ell = \sum_{p \ni \ell} s_p$ involves at most
$2(d{-}1)$ plaquettes.  Changing a neighboring link~$\ell'$
alters the log-density by at most
$2\beta_{\mathrm{eff}} + 2\varepsilon$ per shared plaquette.
The standard Dobrushin bound
(Holley--Stroock comparison~\cite{HolleyStroock}) gives
$C_{\ell\ell'} \leq 1 - e^{-2\beta_{\mathrm{eff}} - 2\varepsilon}$,
and summing over the $6(d{-}1)$ plaquette-neighbors yields
the result.  Since $\varepsilon \to 0$ and
$\beta_{\mathrm{eff}} < \beta_c$ (where Wilson alone satisfies
Dobrushin by Proposition~\ref{prop:dobrushin}), the row sum
is strictly $< 1$ for $\beta$ large enough.
\end{proof}

\subsection{Volume-uniform mixing}

\begin{proposition}\label{prop:mixing}
Under (B1)--(B4) (Theorem~\ref{thm:eac}), the effective measure
$\mu_{k_*}$ satisfies a volume-uniform log-Sobolev inequality:
$\rho_\Lambda(\mu_{k_*}) \geq \rho_* > 0$,
where $\rho_*$ depends only on $G$, $d$, $L$, and~$\beta_c$.
\end{proposition}

\begin{proof}
By Proposition~\ref{prop:perturbed_dob}, the single-link
Dobrushin condition holds.  By
Stroock--Zegarlinski~\cite{StroockZegarlinski},
single-link Dobrushin uniqueness on a compact state space
implies volume-uniform log-Sobolev inequality with
constant $\rho_* = \rho_\ell \cdot (1 - \|C\|_\infty)$,
where $\rho_\ell$ is the single-link conditional LSI constant.
Both $\rho_\ell$ and $\|C\|_\infty$ are volume-independent
by construction.
\end{proof}

\begin{proposition}[Dobrushin analyticity for Schwinger
functions]\label{prop:dob_analytic}
Under (B1)--(B4) (Theorem~\ref{thm:eac}), for any gauge-invariant
local observable~$\mathcal{O}$ on the coarse lattice
$\Lambda_{k_*}$: the connected two-point function at the
Dobrushin scale satisfies
\begin{equation}\label{eq:schwinger_bound}
  \bigl|\langle \mathcal{O}(x)\,\mathcal{O}(y)
  \rangle_{c}^{\mathrm{eff}}
  - \langle \mathcal{O}(x)\,\mathcal{O}(y)
  \rangle_{c}^{W,\beta_c}\bigr|
  \;\leq\; C\,\varepsilon_0\,\|\mathcal{O}\|_\infty^2\,
  e^{-m_0|x-y|/2},
\end{equation}
where $\langle\cdot\rangle_c^{W,\beta_c}$ denotes the
connected correlation in the pure Wilson measure at
coupling~$\beta_c$, $m_0 = -\log(\beta_c/4)$ is the
strong-coupling mass gap, and
$\varepsilon_0 = O(\beta_c^2)$ is the polymer norm
from~\textup{(B4)}.

More generally, for any $n$-point connected correlation:
\begin{equation}\label{eq:schwinger_n}
  \bigl|G_{n,\mathrm{conn}}^{\mathrm{eff}}
  - G_{n,\mathrm{conn}}^{W,\beta_c}\bigr|
  \;\leq\; C_n\,\varepsilon_0\,\|\mathcal{O}\|_\infty^n\,
  e^{-m_0\,\mathrm{diam}(x_1,\ldots,x_n)/2}.
\end{equation}
\end{proposition}

\begin{proof}
The effective Hamiltonian at scale~$k_*$ is
$H_{\mathrm{eff}} = \beta_c\,S_W + R_{k_*}$
with $R_{k_*} = \sum_X a_{k_*}(X)\,\psi_X(U)$
(polymer expansion), where $\psi_X$ is a gauge-invariant
function supported on polymer~$X$ with
$\|\psi_X\|_\infty \leq 1$.

By Proposition~\ref{prop:perturbed_dob}, the measure
$\mu_{\mathrm{eff}}$ satisfies the Dobrushin uniqueness
condition with coefficient
$\alpha = \sum_{\ell'} C_{\ell,\ell'} < 1$.

\smallskip
\noindent\textbf{Analyticity.}
A classical theorem of Dobrushin~\cite{Dobrushin68}
(see also Simon~\cite{Simon_statistical},
Ch.~V.8) states: when the Dobrushin condition holds with
coefficient~$\alpha < 1$, the expectation
$\langle\mathcal{O}\rangle_H$ is analytic in the
Hamiltonian parameters.  Specifically, write
$H(t) = \beta_c S_W + \sum_X t_X\,\psi_X$.
Then $\langle\mathcal{O}\rangle_{H(t)}$ is analytic in
$\{t_X\}$ for $\|t\| < r$, where the analyticity radius
satisfies
\begin{equation}
  r_X \;\geq\; \frac{c\,(1-\alpha)}{|\partial X|}
\end{equation}
and $|\partial X|$ is the number of links in the boundary
of the polymer~$X$.

\smallskip
\noindent\textbf{Convergent expansion.}
Under~(B2) and~(B4):
$|a_{k_*}(X)| \leq \varepsilon_0\,e^{-\gamma\,\diam(X)}$.
For any polymer $X$ with $\diam(X) = R$:
$|\partial X| \leq C\,R^{d-1}$, so
$r_X \geq c'(1-\alpha)/R^{d-1}$.  Meanwhile
$|a_{k_*}(X)| \leq \varepsilon_0\,e^{-\gamma R}$.
For $\varepsilon_0$ small enough (ensured by~(B4)
with $\varepsilon_0 = O(\beta_c^2)$):
\begin{equation}
  |a_{k_*}(X)| < r_X
  \quad \text{for all polymers } X.
\end{equation}
To verify this for \emph{small} polymers ($\diam(X) = R
\leq R_0$ for some fixed $R_0$):
the analyticity radius satisfies $r_X \geq c'(1 - \alpha)$
(since $|\partial X| \leq C\,R_0^{d-1}$ is bounded),
while $|a_{k_*}(X)| \leq \varepsilon_0$.
Here $c' = 1/(2e\,|\partial X|_{\max})$ from
the standard cluster expansion bound
(Simon~\cite{Simon_statistical}, Thm~V.8.3).
For the smallest polymers ($|\partial X| \leq 2d = 8$
in $d = 4$): $c' \geq 1/(16e) \approx 0.023$.
The Dobrushin coefficient at the effective coupling
$\beta_{\mathrm{eff}} \leq \beta_c/L^d$ satisfies
$\alpha \leq 6(d{-}1)(1 - e^{-2\beta_c/L^d - 2\varepsilon_0})$.
For $d = 4$, $L = 2$:
$\beta_c/L^d = 0.0286/16 \approx 0.0018$, so
$\alpha \leq 18 \cdot 0.0044 \approx 0.08$
and $1 - \alpha \geq 0.92$.
Therefore $r_X \geq 0.023 \times 0.92 \approx 0.021$,
while $\varepsilon_0 = O(\beta_c^2) \approx 0.0008 \ll 0.021$.
The condition $\varepsilon_0 < r_X$ is satisfied with
a safety factor of~${\sim}25$.
For \emph{large} polymers ($R > R_0$):
$|a_{k_*}(X)| \leq \varepsilon_0 e^{-\gamma R}$
decays exponentially while $r_X^{-1} \leq
C R^{d-1}/(1-\alpha)$ grows polynomially, so the
bound holds for all $R > R_0(\varepsilon_0)$.
Therefore the Taylor expansion of
$\langle\mathcal{O}\rangle$ in the polymer amplitudes
converges:
\begin{equation}\label{eq:taylor_polymer}
  \langle \mathcal{O}\rangle_{\mathrm{eff}}
  = \langle \mathcal{O} \rangle_{W,\beta_c}
  + \sum_X \frac{\partial\langle\mathcal{O}\rangle}
  {\partial t_X}\bigg|_{t=0} a_{k_*}(X)
  + O(\varepsilon_0^2).
\end{equation}

\smallskip
\noindent\textbf{Decay of the derivative.}
The derivative $\partial\langle\mathcal{O}(x)
\mathcal{O}(y)\rangle_c / \partial t_X|_{t=0}$
is a truncated three-point correlation in the Wilson
measure at~$\beta_c$:
\begin{equation}
  \frac{\partial}{\partial t_X}
  \langle \mathcal{O}(x)\,\mathcal{O}(y)\rangle_c
  \Big|_{t=0}
  = -\langle \mathcal{O}(x);\,\mathcal{O}(y);\,
  \psi_X\rangle_{W,\beta_c},
\end{equation}
which decays exponentially in the minimum distance from
$X$ to $\{x,y\}$ with rate~$m_0$ (strong-coupling
exponential clustering).  Summing over polymers~$X$
with the bound $|a_{k_*}(X)| \leq \varepsilon_0
\,e^{-\gamma\diam(X)}$ gives~\eqref{eq:schwinger_bound}.
The $n$-point bound~\eqref{eq:schwinger_n} follows
identically.
\end{proof}

\begin{remark}[Hypothesis~(H3) resolved]
\label{rem:h3}
Proposition~\ref{prop:dob_analytic} establishes that
the polymer expansion bounds~(B1--B4) imply
$O(\varepsilon_0)$ bounds on individual Schwinger
functions at the Dobrushin scale, with
$\varepsilon_0 = O(\beta_c^2) = O(e^{-c/g^2})$
along the asymptotic freedom trajectory.
This was listed as hypothesis~(H3) in
Proposition~\ref{thm:full}.
Since it now follows from
Proposition~\ref{prop:perturbed_dob}
(Dobrushin uniqueness, conditional on B1--B4)
and the classical Dobrushin analyticity theorem,
(H3) is no longer an independent hypothesis.
\end{remark}

\begin{proposition}[Fine-lattice Schwinger functions from
coarse-lattice bounds]\label{prop:fine_coarse}
Under (B1)--(B4) (Theorem~\ref{thm:eac}), let
$\mathcal{O}$ be a gauge-invariant local observable
on the fine lattice (spacing~$a$).  Then the
connected two-point function satisfies
\begin{equation}\label{eq:fine_coarse}
  \langle \mathcal{O}(x)\,\mathcal{O}(y)\rangle_c^{\text{fine}}
  = \langle \mathcal{O}_{\mathrm{bl}}(x)\,
  \mathcal{O}_{\mathrm{bl}}(y)
  \rangle_c^{\text{coarse}}
  + \delta_{\mathrm{fl}}(x,y),
\end{equation}
where $\mathcal{O}_{\mathrm{bl}}$ is the blocked
observable \textup{(}conditional expectation of
$\mathcal{O}$ given the coarse variables\textup{)},
and the fluctuation contribution satisfies
$|\delta_{\mathrm{fl}}(x,y)| \leq C\,k_*\,\varepsilon_0
\,e^{-\gamma\,|x-y|/(L^{k_*}a)}$.
\end{proposition}

\begin{proof}
Write $\mathcal{O}(x) = \mathcal{O}_{\mathrm{bl}}(x)
+ \delta\mathcal{O}(x)$ where $\delta\mathcal{O}
= \mathcal{O} - E[\mathcal{O}\,|\,U_{\mathrm{coarse}}]$.
By the tower property of conditional expectation:
\begin{align}
  \langle \mathcal{O}(x)\,\mathcal{O}(y)\rangle_c
  &= \langle \mathcal{O}_{\mathrm{bl}}(x)\,
  \mathcal{O}_{\mathrm{bl}}(y)\rangle_c
  + \langle \delta\mathcal{O}(x)\,
  \delta\mathcal{O}(y)\rangle_c,
  \label{eq:tower_decomp}
\end{align}
since the cross terms $\langle
\mathcal{O}_{\mathrm{bl}}(x)\,\delta\mathcal{O}(y)
\rangle_c = 0$ by the definition of conditional
expectation.

The first term in~\eqref{eq:tower_decomp} is controlled
by Proposition~\ref{prop:dob_analytic} (Dobrushin
analyticity at the coarse scale).

For the fluctuation term: $\delta\mathcal{O}(x)$ depends
only on the fine variables within the coarse block
containing~$x$.  When $x$ and $y$ belong to
\emph{different} coarse blocks, the fluctuations
$\delta\mathcal{O}(x)$ and $\delta\mathcal{O}(y)$
are coupled only through the effective action~$R_{k_*}$.
Under (B1)--(B4) \cite{FontesSU2}, the RG integration at each
of the $k_*$ steps produces polymer corrections with
norm~$\leq \varepsilon_0$.  The variance of the fluctuation
decomposes by the tower property:
$\mathcal{O} = \mathcal{O}_{\mathrm{bl}}
+ \sum_{j=1}^{k_*} \delta\mathcal{O}_j$
where $\delta\mathcal{O}_j$ is the fluctuation at
RG step~$j$, satisfying
$\E[\delta\mathcal{O}_j \mid U_{\mathrm{coarse}}^{(j')}]
= 0$ for $j' > j$.  This gives
\begin{equation}
  |\langle \delta\mathcal{O}(x)\,
  \delta\mathcal{O}(y)\rangle_c|
  \;\leq\; \sum_{j=1}^{k_*}
  |\langle \delta\mathcal{O}_j(x)\,
  \delta\mathcal{O}_j(y)\rangle_c|
\end{equation}
(cross terms vanish by the martingale property).
Each individual term is bounded by the polymer expansion
at scale~$j$.  In the worst case, the effective action
at scale~$j$ carries accumulated polymer contributions
from all prior scales, giving
$|\langle \delta\mathcal{O}_j(x)\,
\delta\mathcal{O}_j(y)\rangle_c|
\leq C\,(1+\varepsilon_0)^j\,\varepsilon_0\,
\|\mathcal{O}\|_\infty^2\,
e^{-\gamma|x-y|/(L^j a)}$.
Summing and using $(1+\varepsilon_0)^{k_*}
= \exp(k_*\log(1+\varepsilon_0))
= \exp(O(\varepsilon_0\log\beta))$:
\begin{equation}
  |\langle \delta\mathcal{O}(x)\,
  \delta\mathcal{O}(y)\rangle_c|
  \;\leq\; C\,\beta^{C'\varepsilon_0}\,
  k_*\,\varepsilon_0\,
  \|\mathcal{O}\|_\infty^2\,
  e^{-\gamma\,|x-y|/(L^{k_*}a)}.
\end{equation}
Since $k_* = O(\log\beta)$,
$\varepsilon_0 = O(\beta_c^2)$, and
$\beta^{C'\varepsilon_0} = 1 + O(\varepsilon_0\log\beta)$,
the total correction is
$O(\varepsilon_0\,\log^2\!\beta)$,
which is $o(1)$ relative to the leading coarse-scale
correlation.
\end{proof}

\begin{remark}[Role of polymer bounds in the
fine-coarse transfer]
\label{rem:fine_coarse}
Proposition~\ref{prop:fine_coarse} addresses a
potential concern about the Dobrushin analyticity
approach: since the Dobrushin bound
(Proposition~\ref{prop:dob_analytic}) controls
correlations of \emph{coarse-lattice} variables at
spacing $L^{k_*}a$, one must verify that the
\emph{fine-lattice} Schwinger functions (needed for
the CS equation and OS reconstruction) are similarly
controlled.  The decomposition~\eqref{eq:tower_decomp}
shows that the fine-lattice correlations equal the
coarse-lattice correlations plus a fluctuation term that
is subdominant by virtue of the polymer
bounds~(B1)--(B4).  The key property used is that the RG blocking
kernel factorizes over blocks up to the polymer
corrections---a direct consequence of the polymer
expansion structure~(B1).
\end{remark}

\subsection{Proof of Theorem~\ref{thm:conditional}}

The exponential clustering from Proposition~\ref{prop:mixing}
gives correlation decay on the coarse lattice (spacing
$L^{k_*}\,a$) with rate~$\rho_*$.  For gauge-invariant
observables~$\mathcal{O}_1, \mathcal{O}_2$ at physical
separation~$R$, this translates to
\begin{equation}
  |\langle \mathcal{O}_1(0)\,\mathcal{O}_2(x)\rangle
  - \langle \mathcal{O}_1\rangle\langle\mathcal{O}_2\rangle|
  \;\leq\; C\,e^{-\rho_*\,|x| / (L^{k_*} a)}.
\end{equation}
The mass gap in lattice units is therefore
\begin{equation}
  m \cdot a \;\geq\;
  \frac{\rho_*}{L^{k_*}}
  \;\geq\; \frac{\rho_*}{C_1\,\beta^{1/d}}
  \;=\; \frac{c}{\beta^{1/d}},
\end{equation}
where $c = \rho_* / C_1 > 0$.
For $d = 4$, $G = \SU(2)$:
$\beta_{\mathrm{eff}} \leq \beta_c/L^d = 0.0018$, so
$\rho_* \geq 0.686$
(Proposition~\ref{prop:dobrushin} evaluated at
$\beta_c/L^d$) and $C_1 \leq 9.74$, giving
$c \geq 0.686/9.74 \geq 0.07$. \qed

\begin{remark}[Bare bound versus physical mass gap]
\label{rem:bare_vs_phys}
The bound $m \cdot a \geq c/\beta^{1/d}$ is a \emph{lower}
bound on the \emph{bare} lattice mass gap.
Along the asymptotic freedom trajectory
$\beta \sim (1/g_0^2) \sim \log(1/a)$, this gives a bare
mass $m_{\mathrm{bare}} = (m \cdot a)/a
\geq c/(a\,[\log(1/a)]^{1/d}) \to \infty$.
This divergence is \emph{expected and correct}: in any
non-trivially renormalizable QFT, the bare correlation
functions $G_{\mathrm{bare}}(R;\,a) \to 0$ as $a \to 0$
at fixed physical separation~$R$.  This reflects the
divergence of the wave-function renormalization~$Z(a)$,
not the triviality of the theory.

The \emph{physical} mass gap
$m_{\mathrm{phys}} = \Delta = O(\Lambda_{\mathrm{QCD}})$
is finite and is extracted from the \emph{renormalized}
Schwinger functions
$G_R = Z(a)^n\,G_{\mathrm{bare}}$.
Since $Z(a) \to \infty$ while
$G_{\mathrm{bare}} \to 0$, the product $G_R$ converges
to a nonzero limit---this is precisely the content of the
CS convergence argument (Proposition~\ref{thm:unique}).

Theorem~\ref{thm:obstruction} shows that a
$\beta$-\emph{independent} lower bound
$m \cdot a \geq c_0 > 0$ forces triviality.
The present bound $c/\beta^{1/d} \to 0$ as
$\beta \to \infty$ is $\beta$-dependent and does
\emph{not} trigger this obstruction.  The true mass gap
scales as $m \cdot a \sim e^{-c'\beta}$; our bound
is compatible with (and much weaker than) this scaling.

The role of the mass gap bound in the continuum limit
construction is purely \emph{qualitative}: it
guarantees that $m \cdot a > 0$ for every~$\beta$,
which ensures (i)~absence of IR divergences in the
CS equation, (ii)~exponential clustering for Gibbs
uniqueness, and (iii)~a dominating function for the
RP limit argument (Proposition~\ref{prop:rp}).
The \emph{quantitative} mass gap scaling is determined
by the CS equation and asymptotic freedom, not by the
Dobrushin bound.
\end{remark}


%% ====================================================================
\section{Proof of Theorem~\ref{thm:fsc}:
Finite Size Criterion}\label{sec:fsc}
%% ====================================================================

This is a direct consequence of Theorem~\ref{thm:block_gap}.

\begin{proof}[Proof of Theorem~\ref{thm:fsc}]
The block conditional measure~\eqref{eq:block_cond} on a
$2$-cube $Q = \{0, 1\}^d$ with boundary condition~$\Phi$
is the Wilson measure on $G^{n_{\mathrm{int}}}$ with
$n_{\mathrm{int}} = d \cdot 2^{d-1}$ interior links.

\textbf{Pointwise gap.}
For each $(\beta, \Phi)$: the measure has smooth positive
density on the compact manifold $G^{n_{\mathrm{int}}}$,
so $\gap(\mu_Q(\cdot|\Phi,\beta)) > 0$
(Theorem~\ref{thm:block_gap}(a)).

\textbf{Uniform gap}
(assuming Conjecture~\ref{conj:unique_min}).
By Theorem~\ref{thm:block_gap}(b):
$g(\beta, \Phi) \geq c_1\,\beta$ for $\beta \geq \beta_1$
uniformly in~$\Phi$.  Since $g$ is continuous and
$g \geq c_1\beta$ for $\beta \geq \beta_1$,
the infimum is achieved on the \emph{compact} set
$[0, \beta_1] \times G^{n_{\partial Q}}$, where $g$
is continuous and strictly positive.  Therefore
\begin{equation}
  c_B \;:=\;
  \inf_{\beta \geq 0,\, \Phi \in G^{n_{\partial Q}}}
  g(\beta, \Phi) \;>\; 0.
\end{equation}
This is FSC$(2)$ with constant~$c_B$.

By the Dobrushin--Shlosman complete analyticity
theorem~\cite{DobrushinShlosman,Martinelli1999}:
FSC$(L)$ with $c > 0$ implies unique Gibbs measure
and exponential decay of correlations with rate
$m(c, L, d) > 0$. In particular,
$m(\beta) \cdot a > 0$ for each~$\beta$.
\end{proof}

\begin{remark}[Quantitative status of $c_B$]
\label{rem:fsc_order1}
The compactness argument establishes $c_B > 0$ rigorously.
Numerical evidence from the single-link Witten Laplacian
(Table~\ref{tab:monotone}) suggests $c_B$ may be $O(1)$---the
gap appears monotonically increasing with minimum at $\beta = 0$.
However, the compactness argument alone does not exclude the
possibility that $c_B$ is very small for the full multi-link
block.

For the Dobrushin--Shlosman recursion, what matters is
$c_B > 0$, not its magnitude: the DS theorem gives exponential
clustering with rate depending on~$c_B$, and hence a positive
mass gap.  The \emph{magnitude} of $c_B$ affects only the
quantitative mass gap bound from the single-scale approach.
As discussed in Remark~\ref{rem:triviality}, a
$\beta$-independent bound always gives a trivial continuum limit
regardless of the value of~$c_B$.
\end{remark}

\begin{corollary}[Unconditional results for compact $U(1)$]
\label{cor:U1}
For $G = U(1)$ in $d \geq 2$ dimensions:
\begin{enumerate}[label=\textup{(\alph*)},nosep]
  \item \textup{(Uniform block gap.)}
    The FSC$(2)$ condition holds unconditionally:
    $\gap(\mu_Q(\cdot|\Phi, \beta))
    \geq c_B(d) > 0$ uniformly in $\beta$ and~$\Phi$.
  \item \textup{(Mass gap for $d \leq 3$.)}
    For $d = 2$ and $d = 3$, compact $U(1)$ lattice gauge
    theory is confining for all~$\beta$, and the
    Dobrushin--Shlosman
    theorem~\cite{DobrushinShlosman,Martinelli1999}
    gives unique Gibbs measure and mass gap
    $m(\beta) \cdot a > 0$ for every $\beta > 0$.
\end{enumerate}
For $d = 2$, no assumptions beyond the definition of the
model are required.  For $d = 3$, the proof uses the
G\"opfert--Mack confinement theorem~\cite{GopfertMack82}
(a standard result in constructive gauge theory).
\end{corollary}

\begin{proof}
For $G = U(1)$, Proposition~\ref{prop:unique_abelian} proves
the unique minimum conjecture
(Conjecture~\ref{conj:unique_min}).
Theorem~\ref{thm:block_gap}(b) then gives~(a).

For~(b): the Dobrushin--Shlosman recursion requires that
the uniform block gap propagates to infinite-volume
uniqueness, which needs absence of phase transitions.
In $d = 2$, compact $U(1)$ is exactly solvable via the
transfer matrix, giving uniqueness and mass gap directly.
In $d = 3$, absence of a Coulomb phase (i.e., confinement
for all~$\beta$) was proved by
G\"opfert and Mack~\cite{GopfertMack82}
via a monopole gas representation and Debye screening.
With confinement established, the uniform block gap~(a)
propagates through the DS recursion to give unique
Gibbs measure and $m(\beta) \cdot a > 0$.
\end{proof}

\begin{remark}[The $d = 4$ boundary]
\label{rem:U1_d4}
For $d \geq 4$, compact $U(1)$ lattice gauge theory
exhibits a phase transition between confining and Coulomb
phases~\cite{Guth80}.  The uniform block gap~(a) still holds
(it is a property of the $2$-cube, not of the
infinite-volume theory), but it does not imply mass gap at
weak coupling, where the theory is in the deconfined Coulomb
phase.  This is consistent: the block gap controls local
equilibration but does not override the long-range
correlations responsible for the phase transition.

For non-abelian~$G$ (e.g., $\SU(N)$), no Coulomb phase is
expected in $d = 4$---confinement is believed to hold for
all~$\beta$---so the FSC, if established via the UMC,
would give mass gap at every coupling.  This is one
motivation for resolving Conjecture~\ref{conj:unique_min}
for non-abelian~$G$.
\end{remark}

\begin{remark}[Triviality obstruction]\label{rem:triviality}
Although the FSC constant $c_B$ is $O(1)$, the resulting
mass gap bound is $\beta$-independent:
$m(\beta) \cdot a \geq m_{\mathrm{DS}}(c_B, L, d)$
for all~$\beta$.  This means the correlation length
$\xi = 1/(m \cdot a) \leq C$ remains bounded in lattice
units as $\beta \to \infty$, implying that all connected
correlations vanish in the continuum limit $a \to 0$.
The resulting theory would be \emph{trivial} (free field).

A non-trivial continuum limit requires
$m \cdot a \to 0$ as $\beta \to \infty$---specifically,
$m \cdot a \sim \Lambda_{\mathrm{QCD}} \cdot a
\sim e^{-c_{\mathrm{AF}} \beta}$ from asymptotic freedom.
This $\beta$-dependence can only be captured by a
multi-scale argument that tracks the effective coupling
through the RG.  Theorem~\ref{thm:conditional} achieves
this with the polynomial bound $m \cdot a \geq c/\beta^{1/d}$,
which allows $\xi \sim \beta^{1/d} \to \infty$.

The FSC result (Theorem~\ref{thm:fsc}) is therefore
\emph{qualitatively correct} (mass gap exists at every
coupling) but \emph{quantitatively insufficient} for a non-trivial
continuum limit.  The multi-scale analysis of
Theorem~\ref{thm:conditional}, with the effective action
control~\cite{FontesSU2}, is essential for this purpose.
\end{remark}

\begin{lemma}[High-temperature Dobrushin]\label{lem:high_temp_dob}
For $\beta \leq \beta_0 = N/(4 \cdot 6(d{-}1))$, the lattice
Yang--Mills theory satisfies single-link Dobrushin uniqueness
with $\sum_{\ell'} C_{\ell\ell'} \leq 6(d{-}1) \cdot \beta
\cdot 4/N < 1$.
\end{lemma}


%% ====================================================================
\section{Proof of Theorem~\ref{thm:obstruction}:
The Triviality Obstruction}\label{sec:obstruction}
%% ====================================================================

\subsection{Proof of Theorem~\ref{thm:obstruction}(a): Triviality}

If $m(\beta) \cdot a \geq c_0 > 0$ uniformly in~$\beta$, then
the correlation length $\xi = 1/(m \cdot a) \leq 1/c_0$
in lattice units.  For gauge-invariant observables
$\mathcal{O}_1, \mathcal{O}_2$ at separation~$r$ lattice spacings:
\begin{equation}
  |\langle \mathcal{O}_1(0)\,\mathcal{O}_2(x)\rangle_{\mathrm{conn}}|
  \;\leq\; C\,e^{-c_0\,r}.
\end{equation}
At physical separation $R = r \cdot a$ with $r = R/a$:
\begin{equation}
  |\langle \mathcal{O}_1(0)\,\mathcal{O}_2(R)\rangle_{\mathrm{conn}}|
  \;\leq\; C\,e^{-c_0\,R/a} \;\to\; 0 \quad \text{as } a \to 0
\end{equation}
for any fixed $R > 0$.  All connected Schwinger functions
vanish in the continuum limit.  The resulting theory
satisfies clustering with infinite rate, which by the
Jost--Schroer theorem implies a free (generalized free) field.
\qed

\subsection{What the entropy framework provides}

At the RG scale $k_*$ where $\beta_{\mathrm{eff}} < \beta_c$,
the entropy bound (Theorem~\ref{thm:entropy}) gives
\begin{equation}
  \KL(\mu_{k_*} \| \nu_{k_*})
  \;\leq\; \frac{C}{\beta^2}.
\end{equation}
By Pinsker's inequality:
\begin{equation}\label{eq:pinsker}
  \|\mu_{k_*} - \nu_{k_*}\|_{\TV}
  \;\leq\; \sqrt{C / (2\beta^{2})}
  \;\leq\; C' / \beta.
\end{equation}
\textup{(}If Conjecture~\ref{conj:entropy_geom}
holds, the sharper bound
$\KL \leq C/\beta^{2+\alpha}$ with
$\alpha > 0$ follows, but the uniform
bound suffices for all results below.\textup{)}

\subsection{What is still needed}

The Dobrushin condition (Proposition~\ref{prop:perturbed_dob})
requires \emph{pointwise} control: for each link~$\ell$,
the oscillation $\osc_\ell(R_{k_*}) \leq \varepsilon$
where $R_{k_*}$ is the non-Wilson part of the effective action.

The entropy bound~\eqref{eq:pinsker} gives
\emph{average} (TV) closeness to the Gaussian, which
implies closeness to Haar at small $\beta_{\mathrm{eff}}$.
However, TV closeness does not directly bound the
per-link oscillation of the log-density ratio.

The obstruction is localized: at weak coupling
($\beta_{\mathrm{eff}} \gg 1$), the per-plaquette bound
ensures that non-Gaussian corrections are $O(1/\beta^2)$
per plaquette, giving excellent pointwise control.
At strong coupling ($\beta_{\mathrm{eff}} \ll 1$), the
cluster expansion converges and gives pointwise bounds.
The gap is at \emph{intermediate coupling}
($\beta_{\mathrm{eff}} \approx 1$): here the fields are
$O(1)$, the quartic approximation fails, and the cluster
expansion is not yet convergent.

\subsection{The intermediate coupling regime}

The RG passes through intermediate coupling in
$O(1)$ steps (specifically, $\lceil \log(1/\beta_c) /
(d\log L)\rceil \leq 3$ steps for $d = 4$, $L = 2$).
During these steps, the effective action accumulates
corrections of $O(1)$ per link from the integration over
interior variables at coupling $\beta_{\mathrm{eff}} \approx 1$.

The companion paper~\cite{FontesSU2} provides
the required pointwise control through this regime for
$G = \SU(2)$,
establishing that $R_{k_*}$ has polymer structure with
$\osc_\ell(R_{k_*}) \leq \varepsilon(\beta) \to 0$
as $\beta \to \infty$.
(Ba{\l}aban's original
work~\cite{Balaban85,Balaban89} contains analogous
results in greater generality.)

\begin{proposition}[Sufficient condition for unconditional
mass gap]\label{prop:sufficient}
If the effective measure at any RG scale $k_*$ where
$\beta_{\mathrm{eff}}(k_*) < \beta_c/2$ can be written as
\begin{equation}
  d\mu_{k_*}(x) \propto \exp\!\bigl(-\beta_{\mathrm{eff}}
  \,S_W(x) - R_{k_*}(x)\bigr)\,\prod_\ell dU_\ell
\end{equation}
with $\sum_{X \ni \ell} \sup |a_{k_*}(X)|
\leq \varepsilon < \varepsilon_0(d)$ where
\begin{equation}\label{eq:eps0}
  \varepsilon_0 = \frac{1}{2}\Bigl(
  -\log\Bigl(1 - \frac{1}{6(d{-}1)}\Bigr) - \beta_c\Bigr),
\end{equation}
then $m \cdot a \geq c / \beta^{1/d}$ unconditionally.
For $d = 4$: $\varepsilon_0 \approx 0.030$.
\end{proposition}

\begin{proof}
The perturbed Dobrushin row sum becomes
$6(d{-}1)(1 - e^{-2(\beta_c/2 + \varepsilon)})
\leq 6(d{-}1)(1 - e^{-\beta_c - 2\varepsilon})$.
By~\eqref{eq:eps0}:
$\beta_c + 2\varepsilon < \beta_c + 2\varepsilon_0
= -\log(1 - 1/(6(d{-}1)))$,
so the row sum is $< 1$.  The rest follows as in
\S\ref{sec:conditional}.
\end{proof}


%% ====================================================================
\section{Continuum Limit and Uniqueness}\label{sec:continuum}
%% ====================================================================

\subsection{Lattice mass gap and continuum}

Under Theorem~\ref{thm:conditional}, the lattice mass gap
satisfies $m \cdot a \geq c / \beta^{1/4}$ for $d = 4$.
Along the asymptotic freedom trajectory
$\beta(a) \sim c_{\mathrm{AF}} \log(1/a)$, the
\emph{bare} mass extracted from the bound diverges:
\begin{equation}
  m_{\mathrm{bare\text{-}bound}}(a)
  = \frac{m \cdot a}{a}
  \;\geq\; \frac{c}{a\,[\log(1/a)]^{1/4}}
  \;\to\; \infty \quad \text{as } a \to 0.
\end{equation}
This divergence is \emph{expected}, not pathological:
it reflects the fact that the bound $c/\beta^{1/d}$
is a loose lower bound on $m \cdot a$, far above the
true scaling $m \cdot a \sim e^{-c'\beta}$.
The \emph{physical} mass gap
$m_{\mathrm{phys}} = O(\Lambda_{\mathrm{QCD}})$
is finite and is extracted from the renormalized
Schwinger functions via the CS equation, not from
the Dobrushin bound (see Remark~\ref{rem:bare_vs_phys}).

\subsection{Subsequential continuum limit}

For gauge-invariant observables: the lattice states
$\{\omega_a\}$ form a tight family (Wilson loops are
uniformly bounded by~1 and exponential clustering is
uniform in~$a$).  By Prokhorov's theorem, a
subsequential limit $\omega_*$ exists.

The limit inherits:
\begin{enumerate}[label=(\roman*),nosep]
  \item Reflection positivity (from the lattice, preserved
    under weak limits).
  \item $\SO(4)$ covariance (from the lattice symmetry
    group in the continuum limit).
  \item Exponential clustering with rate $m_* > 0$
    (from the uniform lattice bound).
\end{enumerate}
By Osterwalder--Schrader
reconstruction~\cite{OsterwalderSchrader}:
$\omega_*$ yields a relativistic QFT satisfying
the Wightman axioms with mass gap $\Delta = m_* > 0$.

\subsection{Reduction to the two-point function}

We now address uniqueness of the continuum limit.  The
strategy reduces uniqueness of \emph{all} Schwinger
functions to convergence of the \emph{two-point function}.

\begin{proposition}[SD uniqueness in the massive regime]
\label{prop:sd_unique}
Let $\{S_n\}_{n \geq 2}$ and $\{S'_n\}_{n \geq 2}$ be
two systems of connected Schwinger functions for $\SU(N)$
Yang--Mills theory, each satisfying the
Osterwalder--Schrader axioms, the Schwinger--Dyson equations,
and exponential clustering with mass gap $\Delta > 0$.
If $S_2 = S'_2$
\textup{(}the two-point functions agree\textup{)},
then $S_n = S'_n$ for all~$n$.
\end{proposition}

\begin{proof}
\textbf{Step 1: Structure of the SD equations.}
The Schwinger--Dyson equations for gauge-invariant
Schwinger functions of a local operator~$\mathcal{O}$
(of canonical dimension~$d_{\mathcal{O}}$) express
$S_{n+1}$ in terms of lower correlators by integrating
the equations of motion inside the functional integral.
Schematically:
\begin{equation}\label{eq:sd_structure}
  S_{n+1}(x_1,\ldots,x_{n+1}) =
  \sum_{\substack{\text{tree-level}}}
  V \cdot \bigl[(G * S_k)(x_I)\bigr]
  \cdot S_{n+1-k}(x_{I^c})
  + \mathcal{R}_{n+1}[S_2; S_{n+1}],
\end{equation}
where $G = S_2$ is the propagator (two-point function),
$V$ denotes vertex factors depending on the coupling
$g_R$ and the operator~$\mathcal{O}$, the tree-level sum
runs over partitions $I \sqcup I^c = \{1,\ldots,n+1\}$
with $|I| = k$ and involves only $S_j$ with $j \leq n$,
and the remainder
$\mathcal{R}_{n+1}[S_2; S_{n+1}]$ arises from
one-loop and higher corrections involving $S_{n+1}$ itself.

\smallskip
\textbf{Step 2: Contraction of the self-consistent term.}
The one-loop remainder has the form
\[
  \mathcal{R}_{n+1}[G; S_{n+1}](x_1,\ldots,x_{n+1})
  = \int V_{\mathrm{loop}}\,G(z{-}z')\,
  K(x_1,\ldots,x_{n+1}; z, z')\,dz\,dz',
\]
where $K$ involves $S_{n+1}$ evaluated at shifted
arguments.  The operator norm satisfies
\begin{equation}\label{eq:loop_bound}
  \|\mathcal{R}_{n+1}[G; \cdot]\|_{\mathrm{op}}
  \;\leq\; C_n\,g_R^2\,\|G\|_{L^1}
  \;\leq\; C_n\,g_R^2 / \Delta^{d-2},
\end{equation}
since $|G(x)| \leq C\,e^{-\Delta|x|}$ gives
$\|G\|_{L^1} \leq C'/\Delta^{d-2}$.  Here $C_n$
depends polynomially on~$n$ (from the number of
contractions).

For $d = 4$: $\|G\|_{L^1} \leq C/\Delta^2$, and the
dimensionless control parameter is $g_R^2/\Delta^2$.
Along the asymptotic freedom trajectory,
$g_R^2 \to 0$ and $\Delta = O(\Lambda_{\mathrm{QCD}})$
is fixed, so $g_R^2/\Delta^2 \to 0$.

\smallskip
\textbf{Step 3: Induction on $n$.}
For $g_R^2\,\|G\|_{L^1}$ sufficiently small
(achieved along the AF trajectory for all
$a \leq a_0$), the operator
$I - \mathcal{R}_{n+1}[G; \cdot]$ is invertible
by \eqref{eq:loop_bound}, and
\begin{equation}\label{eq:sd_solved}
  S_{n+1} = (I - \mathcal{R}_{n+1}[G;\cdot])^{-1}
  \Bigl(\sum_{\text{tree}} V \cdot
  (G * S_k) \cdot S_{n+1-k}\Bigr).
\end{equation}
The right-hand side depends only on $G = S_2$ and
$\{S_j\}_{j \leq n}$.  Proceed by induction:

\noindent\textsc{Base}: $S_2 = S'_2$ by hypothesis.

\noindent\textsc{Step}: Suppose $S_j = S'_j$ for all
$j \leq n$.  Then the right-hand side
of~\eqref{eq:sd_solved} is identical for
$S_{n+1}$ and $S'_{n+1}$, since all inputs agree.
Therefore $S_{n+1} = S'_{n+1}$.

\smallskip
\textbf{Step 4: Growth control via asymptotic freedom.}
The Neumann inverse $(I - \mathcal{R})^{-1}$ introduces
a factor $(1 - C_n g_R^2/\Delta^{d-2})^{-1}$ at each
induction step.  For the induction to close at all~$n$:
we need $C_n g_R^2/\Delta^{d-2} < 1$ for every~$n$.
Since $C_n = O(n^P)$ for some power~$P$ (from vertex
combinatorics), this fails for $n > n_{\max}(g_R)$
at any fixed $g_R > 0$.

\emph{Resolution.}
The SD equations for the \emph{lattice} Schwinger
functions $S_n^{(a)}(x_1, \ldots, x_n)$ at fixed
physical separations $|x_i - x_j| \sim 1/\Lambda_{\mathrm{QCD}}$
involve loop integrals where internal momenta range
up to $\pi/a$ (the UV cutoff).
By asymptotic freedom, the effective coupling
controlling these loops is the \emph{bare} coupling
$g^2(a) = 1/(2\beta_0\log(1/a)) \to 0$ as $a \to 0$,
not the physical coupling $g_R(\mu)$ at the
external momentum scale.
Therefore the contraction
bound~\eqref{eq:loop_bound} involves $g^2(a)$,
not $g_R(\Lambda_{\mathrm{QCD}})$:
\[
  \|\mathcal{R}_{n+1}\|_{\mathrm{op}}
  \;\leq\; C_n\,g^2(a) / \Delta^{d-2}.
\]
For any fixed~$N$, there exists $a_0(N)$ such that
this is $< 1$ for all $n \leq N$ and all $a < a_0(N)$.
In the continuum limit $a \to 0$, the bound holds
for all~$n$ simultaneously, so the lattice SD
equations uniquely determine $S_n^{(a)}$ from $S_2^{(a)}$
for each~$a$.

The passage to the continuum ($a \to 0$) then uses
the CS equation (Lemma~\ref{lem:irrelevant} and
Proposition~\ref{thm:unique}), which shows that
the lattice Schwinger functions converge to a unique
limit.  The uniqueness chain is:
$S_2^{(a)}$ (given for each~$a$)
$\xrightarrow{\text{SD on lattice}}$
all $S_n^{(a)}$ (unique for each~$a$)
$\xrightarrow{\text{CS convergence}}$
all $S_n^{(\mathrm{cont})}$ (unique limit).
\end{proof}

\subsection{Callan--Symanzik equation on the lattice}

The Callan--Symanzik (CS) equation governs how renormalized
Green's functions evolve with the lattice spacing.  For the
renormalized $n$-point function of a gauge-invariant local
operator of dimension~$\Delta_{\mathcal{O}}$:
\begin{equation}\label{eq:cs_lattice}
  \Bigl[a\frac{\partial}{\partial a}
  + \beta_L(g_L, a)\frac{\partial}{\partial g_L}
  + n\,\gamma_L(g_L, a)\Bigr]
  G_R^{(n)}(p_i; g_L, a) = 0,
\end{equation}
where $\beta_L$ and $\gamma_L$ are the \emph{exact} lattice
beta-function and anomalous dimension.

\begin{lemma}[Irrelevant corrections]\label{lem:irrelevant}
Assume (B1)--(B4) \cite{FontesSU2} and the mass gap
\textup{(}Theorem~\ref{thm:conditional}\textup{)}.
The lattice beta-function and anomalous dimension satisfy
\begin{equation}\label{eq:irrelevant}
  \beta_L(g, a) = \beta(g) + O(a^2 g^2),
  \qquad
  \gamma_L(g, a) = \gamma(g) + O(a^2 g^2),
\end{equation}
where $\beta(g)$ and $\gamma(g)$ are the universal
\textup{(}scheme-independent\textup{)} continuum functions,
and the $O(a^2 g^2)$ corrections are bounded uniformly
in~$\beta$.
\end{lemma}

\begin{proof}
The proof combines a perturbative result from the
literature with the non-perturbative control established
in Proposition~\ref{prop:dob_analytic}.

\smallskip
\noindent\textbf{Perturbative part.}
The Wilson action satisfies the hypotheses of Reisz's
lattice power-counting theorem~\cite{Reisz1988}:
\begin{enumerate}[label=(R\arabic*),nosep]
  \item \emph{Correct continuum limit.}
    With $U_\ell = e^{iaA_\mu}$, the plaquette action
    expands as $(1{-}(1/N)\Re\tr U_p)
    = (a^4/2N)\tr F_{\mu\nu}^2 + O(a^6)$,
    so $S_W = S_{\mathrm{YM}} + O(a^2)$.
  \item \emph{Propagator bound.}
    The lattice gluon propagator satisfies
    $G_L(p)^{-1} = (4/a^2)\sum_\mu \sin^2(p_\mu a/2)
    \cdot(\text{gauge factor})
    = |p|^2(1 + O(a^2|p|^2))$.
  \item \emph{Vertex bounds.}
    The $3$- and $4$-gluon lattice vertices match the
    continuum vertices to $O(a^2)$, by expansion of the
    Wilson action to cubic and quartic order.
  \item \emph{Renormalizability.}
    YM in $d = 4$ is perturbatively renormalizable.
\end{enumerate}
By Reisz's theorem~\cite{Reisz1988}
(extended to gauge theories in~\cite{Reisz89}):
every lattice Feynman integral contributing to
$\beta_L$ or $\gamma_L$ at $K$ loops matches its
continuum counterpart up to $O(a^2)$ corrections
(from dimension-$6$ irrelevant operators).
The first two coefficients $\beta_0, \beta_1$ of
$\beta(g) = -\beta_0 g^3 - \beta_1 g^5 - \cdots$ are
universal (scheme-independent) and match
exactly~\cite{Luscher_Weisz}.
Therefore, to any finite perturbative order~$K$:
\begin{equation}\label{eq:pert_matching}
  |\beta_L^{(K)}(g) - \beta^{(K)}(g)|
  \;\leq\; C_K\,a^2\,g^{2K+4}\,[\log(1/a)]^{P_K},
\end{equation}
and similarly for $\gamma_L$.

\smallskip
\noindent\textbf{Non-perturbative part.}
The correction beyond $K$-th order perturbation theory
has two sources:
\begin{enumerate}[label=(\roman*),nosep]
  \item The perturbative remainder of order $g^{2K+2}$,
    bounded by $C_K g^{2K+2}$.
  \item The genuinely non-perturbative contribution
    (instanton-type).  By
    Proposition~\ref{prop:dob_analytic} (Dobrushin
    analyticity), applied to the effective theory at the
    Dobrushin scale, this is bounded by
    $C\,\varepsilon_0 = O(\beta_c^2)$.
\end{enumerate}
Along the asymptotic freedom trajectory
$g^2(a) = 1/(2\beta_0\log(1/a))$:
\begin{equation}
  \varepsilon_0 = O(\beta_c^2) = O\bigl(\beta/L^{k_* d}\bigr)^2
  = O(e^{-c/g^2}) = O(a^{2c\beta_0}).
\end{equation}
For $\SU(2)$ with $c = 8\pi^2/N$:
$2c\beta_0 \approx 3.67$, so the non-perturbative
correction is $O(a^{3.67}) = o(a^2)$.

The mass gap ensures IR contributions are bounded by
$O(e^{-m \cdot a}) = o(a^k)$ for all~$k$.

\smallskip
\noindent\textbf{Combined.}
The total error $|\beta_L(g,a) - \beta(g)|$ is bounded by
three terms:
\begin{enumerate}[label=(\alph*),nosep]
  \item \emph{Perturbative lattice artifacts}
    (from~\eqref{eq:pert_matching}):
    $C_K\,a^2\,g^{2K+4}\,[\log(1/a)]^{P_K}$.
  \item \emph{Perturbative remainder} at order $K$:
    $C_K\,g^{2K+2}$.
  \item \emph{Non-perturbative (instanton) contribution}:
    $C'\,a^{2c\beta_0}$ with $2c\beta_0 \approx 3.67$
    for $\SU(2)$.
\end{enumerate}
Along the asymptotic freedom trajectory
$g^2 = 1/(2\beta_0\log(1/a))$, each term is absorbed
into $O(a^2 g^2)$ as follows.

For~(b): $g^{2K+2} = [2\beta_0\log(1/a)]^{-(K+1)}$.
This is $o(a^2 g^2) = o(a^2/\log(1/a))$ provided
$(K+1)\log\log(1/a) \gg 2\log(1/a)$, which holds for
$K \geq C/\log\log(1/a)$.  Choose $K = K(a)$ growing
with $\log\log(1/a)$; then (b)~$= o(a^2 g^2)$.

For~(c): $a^{3.67} / (a^2 g^2) = a^{1.67}\log(1/a) \to 0$
as $a \to 0$, so (c)~$= o(a^2 g^2)$.

For~(a): with the above choice of~$K$,
$a^2 g^{2K+4}[\log(1/a)]^{P_K}
= a^2 g^4 \cdot g^{2K}[\log(1/a)]^{P_K}$,
and $g^{2K} = o(1)$ since $K \to \infty$, giving
(a)~$= o(a^2 g^4) \subseteq O(a^2 g^2)$.

Combining:
\begin{equation}
  |\beta_L(g,a) - \beta(g)|
  \;\leq\; C\,a^2\,g^2
  + o(a^2 g^2)
  + O(a^{3.67})
  \;\leq\; C''\,a^2\,g^2,
\end{equation}
establishing~\eqref{eq:irrelevant}.  The key point is that
the non-perturbative error~(c), controlled by Dobrushin
analyticity at the effective coupling, decays \emph{faster}
than $a^2$ because $2c\beta_0 > 2$; and the perturbative
remainder~(b) is tamed by choosing the expansion order~$K$
to grow logarithmically with $1/a$.
\end{proof}

\begin{remark}[Hypothesis~(H1) resolved]
\label{rem:h1}
Lemma~\ref{lem:irrelevant} establishes the $O(a^2)$
irrelevant corrections to the lattice CS equation, which
was listed as hypothesis~(H1) in
Proposition~\ref{thm:full}.  The proof combines Reisz's
power-counting theorem~\cite{Reisz1988,Reisz89}
(a standard result in the lattice QFT literature)
with the Dobrushin analyticity
(Proposition~\ref{prop:dob_analytic}) for the
non-perturbative part.  Since both ingredients are
available, (H1) is no longer an independent hypothesis.
\end{remark}

\subsection{Convergence via Gronwall's inequality}

\begin{proposition}[Euclidean continuum limit]
\label{thm:unique}
Assume (B1)--(B4) \cite{FontesSU2}.
Then:
\begin{enumerate}[label=\textup{(\alph*)},nosep]
  \item The mass gap satisfies
    $m \cdot a \geq c/\beta^{1/d}$
    \textup{(}Theorem~\ref{thm:conditional}\textup{)}.
  \item The renormalized Schwinger functions of
    gauge-invariant local operators converge as
    $a \to 0$ along the asymptotic freedom trajectory.
  \item The Euclidean continuum limit is unique, determined
    by the gauge group~$G$ and the dynamical scale
    $\Lambda_{\mathrm{QCD}}$.
  \item The theory is interacting: $\lambda_4 \neq 0$
    \textup{(}Proposition~\ref{thm:nontrivial}\textup{)}.
\end{enumerate}
Furthermore:
\begin{enumerate}[label=\textup{(\alph*)},nosep]
  \setcounter{enumi}{4}
  \item The continuum Schwinger functions satisfy reflection
    positivity \textup{(}inherited from the lattice Wilson
    action via dominated convergence,
    Proposition~\ref{prop:rp}\textup{)}.
  \item By Osterwalder--Schrader
    reconstruction, the theory defines a relativistic
    QFT satisfying the Wightman axioms with mass gap
    $\Delta > 0$.
\end{enumerate}
\end{proposition}

\begin{proof}
\textbf{Step 1: Running coupling convergence.}
Let $g_L(a)$ denote the lattice running coupling,
satisfying $a\,dg_L/da = \beta_L(g_L, a)$, and let
$\bar{g}(a)$ denote the continuum running coupling,
satisfying $a\,d\bar{g}/da = \beta(\bar{g})$, with the
same initial condition at a reference scale~$a_0$.
Define $\delta g(a) = g_L(a) - \bar{g}(a)$.  Then
\begin{equation}\label{eq:delta_g}
  a\,\frac{d(\delta g)}{da}
  = \underbrace{[\beta(g_L) - \beta(\bar{g})]}_{
  = \beta'(\xi)\,\delta g}
  + \underbrace{[\beta_L(g_L, a) - \beta(g_L)]}_{
  = O(a^2 g_L^2)},
\end{equation}
where $\xi$ lies between $g_L$ and $\bar{g}$
(mean value theorem) and the $O(a^2 g_L^2)$ bound
is Lemma~\ref{lem:irrelevant}.

Since $\beta'(g) = -3\beta_0 g^2 + O(g^4)$
is \emph{negative} for $g$ small (asymptotic freedom),
$\beta'(\xi) \leq -2\beta_0 \xi^2 < 0$ for
$\xi \leq g_0$ with $g_0$ sufficiently small.
This makes the homogeneous part of~\eqref{eq:delta_g}
\emph{contractive}: perturbations in $g$ are damped.

Define the integrating factor
$\Phi(a, a') :=
\exp(\int_a^{a'} \beta'(\xi(s))\,ds/s) \leq 1$
(since $\beta' < 0$).  By the variation-of-constants
formula applied to~\eqref{eq:delta_g}:
\begin{equation}\label{eq:gronwall_bound}
  |\delta g(a)|
  \;\leq\; \int_a^{a_0} \Phi(a, a')\,
  C\,(a')^2\,g_L(a')^2\,\frac{da'}{a'}
  \;\leq\; C \int_a^{a_0}
  (a')^2\,g_L(a')^2\,\frac{da'}{a'}.
\end{equation}
Along the asymptotic freedom trajectory,
$g_L(a)^2 \sim 1/(2\beta_0\log(1/a))$, so
\begin{equation}
  |\delta g(a)|
  \;\leq\; \frac{C}{(2\beta_0)^2}
  \int_a^{a_0}
  \frac{a'\,da'}{[\log(1/a')]^2}
  \;\leq\; \frac{C\,a_0^2}{[\log(1/a_0)]^2}
  \;\to\; 0 \quad \text{as } a_0 \to 0.
\end{equation}
Here the integral converges because the
$a'\,da' = O((a')^2)$ factor overwhelms the
logarithmic denominator.  The lattice running
coupling converges uniformly to the continuum
running coupling on $(0, a_0]$.

\smallskip
\textbf{Step 2: Green's function convergence.}
The solution of~\eqref{eq:cs_lattice} along
the characteristic curves is
\begin{equation}\label{eq:cs_solution}
  G_R^{(n)}(p; a)
  = G_R^{(n)}(p; a_0)\,
  \exp\!\Bigl(-n\int_a^{a_0}
  \gamma_L(g_L(a'), a')\,\frac{da'}{a'}\Bigr).
\end{equation}
Split the exponent into three terms:
\begin{align}
  \int_a^{a_0} \gamma_L(g_L, a')\,\frac{da'}{a'}
  &= \underbrace{\int_a^{a_0}
  \gamma(\bar{g}(a'))\,\frac{da'}{a'}}_{=: I_1(a)}
  \;+\; \underbrace{\int_a^{a_0}
  [\gamma_L(g_L, a') - \gamma(g_L)]\,
  \frac{da'}{a'}}_{=: I_2(a)} \notag \\
  &\quad +\; \underbrace{\int_a^{a_0}
  [\gamma(g_L(a')) - \gamma(\bar{g}(a'))]\,
  \frac{da'}{a'}}_{=: I_3(a)}.
  \label{eq:exponent_split}
\end{align}

\noindent$I_1(a)$: Since
$\gamma(\bar{g}) = \gamma_0 \bar{g}^2 + O(\bar{g}^4)$
with $\bar{g}(a)^2 \sim 1/(2\beta_0\log(1/a))$,
the integrand is $O(1/\log(1/a'))$, which is integrable
on $(0, a_0]$.  Hence $I_1(a) \to
\int_0^{a_0} \gamma(\bar{g})\,da'/a'$, a finite limit.

\noindent$I_2(a)$: By Lemma~\ref{lem:irrelevant},
$|\gamma_L(g,a) - \gamma(g)| \leq C\,a^2\,g^2$.
So $|I_2(a)| \leq C \int_0^{a_0}
a'\,g_L^2(a')\,da'/a' = O(a_0^2/[\log(1/a_0)]^2)$,
which is finite and independent of~$a$.

\noindent$I_3(a)$: Since $\gamma$ is Lipschitz near $g = 0$
with $|\gamma(g) - \gamma(g')| \leq C|g - g'|\,(|g| + |g'|)$,
we have $|I_3(a)| \leq C\,\|\delta g\|_\infty \cdot
\int_0^{a_0} g_L(a')\,da'/a'$.  By Step~1,
$\|\delta g\|_\infty \to 0$, so $I_3(a) \to 0$.

All three terms converge.  Therefore the exponential
factor in~\eqref{eq:cs_solution} has a definite limit
$Z_n := \lim_{a \to 0} \exp(-n\int_a^{a_0}\gamma_L\,da'/a')
\in (0,\infty)$,
and $G_R^{(n)}(p; a) \to G_R^{(n)}(p; a_0) \cdot Z_n$
as $a \to 0$.

\smallskip
\textbf{Step 3: Unique initial condition.}
By Proposition~\ref{prop:perturbed_dob} (Dobrushin uniqueness
at the Dobrushin scale), the lattice Yang--Mills
measure at any fixed~$\beta$ in the weak coupling regime
has a unique infinite-volume Gibbs measure.
(Dobrushin uniqueness of the effective measure at
scale~$k_*$ implies uniqueness of the original
measure, since the conditional distributions within
blocks are determined by the effective state.)
Therefore $G_R^{(n)}(p; a_0)$ is a \emph{uniquely determined}
function of $\beta(a_0)$---there is no subsequence ambiguity
at fixed lattice spacing.

\smallskip
\textbf{Step 4: Independence of reference scale.}
Changing $a_0 \to a_0'$: the CS
equation~\eqref{eq:cs_solution} connects the Green's
functions at different reference scales by the
multiplicative factor
$\exp(-n\int_{a_0}^{a_0'}\gamma_L\,da'/a')$, which is
uniquely determined by the lattice RG flow.
The limit as $a \to 0$ is
independent of the choice of~$a_0$.

\smallskip
\textbf{Step 5: All Schwinger functions.}
Steps~1--4 establish convergence of the two-point
function $G_R^{(2)}(p; a) \to G_R^{(2)}(p)$ as $a \to 0$.
By Proposition~\ref{prop:sd_unique} (SD uniqueness),
the two-point function determines all higher
$G_R^{(n)}$ uniquely via the massive SD equations.
Hence all Schwinger functions converge.

\smallskip
\textbf{Step 6: Wightman axioms.}
The convergent Schwinger functions satisfy the OS axioms:
\begin{enumerate}[label=(\roman*),nosep]
  \item Euclidean covariance: inherited from the lattice
    symmetries, which converge to continuous rotations.
  \item Reflection positivity: Proposition~\ref{prop:rp}.
  \item Exponential clustering with gap
    $\Delta \geq \liminf_{a \to 0} m(a) \cdot a /a > 0$,
    inherited from the lattice mass gap.
\end{enumerate}
OS reconstruction~\cite{OsterwalderSchrader} yields a
Wightman QFT with mass gap $\Delta > 0$.

\smallskip
\textbf{Step 7: Dimensional transmutation.}
The continuum CS equation $[a\partial_a + \beta(g)\partial_g
+ n\gamma(g)]G^{(n)} = 0$ has a one-parameter family of
solutions, parametrized by the RG-invariant scale
$\Lambda_{\mathrm{QCD}}$.  The lattice defines
$\Lambda_{\mathrm{QCD}}$ via
\begin{equation}
  \Lambda_{\mathrm{QCD}}
  = \frac{1}{a}\,e^{-1/(2\beta_0 g^2(a))}
  \,(2\beta_0 g^2(a))^{-\beta_1/(2\beta_0^2)}\,
  [1 + O(g^2)],
\end{equation}
which converges as $a \to 0$ (since $\delta g \to 0$
by Step~1).
The theory is uniquely determined by~$G$
and~$\Lambda_{\mathrm{QCD}}$.
\end{proof}

\begin{remark}[Role of each assumption]
\label{rem:uniqueness_roles}
The proof uses three distinct inputs:

\smallskip\noindent
(i) \textbf{(B1)--(B4) \cite{FontesSU2}:}
Bounds the non-perturbative contributions to the lattice
CS coefficients.  Without these bounds, the irrelevant
corrections in~\eqref{eq:irrelevant} could diverge at
intermediate coupling, making the Gronwall integral
diverge.

\smallskip\noindent
(ii) \textbf{Mass gap~(B5):} Provides the IR cutoff
ensuring that all sums over lattice sites converge and
the SD iteration contracts.  Without the mass gap, the
CS equation could have IR-divergent corrections.

\smallskip\noindent
(iii) \textbf{Gibbs uniqueness
(Proposition~\ref{prop:perturbed_dob}):} Gives the
unique initial condition at the reference scale~$a_0$.
Without Gibbs uniqueness, the Green's functions at
fixed~$a_0$ could be multi-valued (depending on boundary
conditions or the thermodynamic limit procedure),
destroying the deterministic CS evolution.
This follows from Dobrushin uniqueness at the Dobrushin
scale (conditional on B1--B4), independently of the
unconditional FSC~(Theorem~\ref{thm:fsc},
conditional on the unique minimum conjecture).

\smallskip\noindent
Asymptotic freedom is also essential: without $g(a) \to 0$,
the irrelevant corrections $O(a^2 g^2)$ might fail to be
integrable.  This is consistent with the expected triviality
of QED$_4$, where $g(a) \to \infty$ and the continuum
limit does not exist.
\end{remark}

\begin{remark}[No Borel summability needed]
\label{rem:no_borel}
A common approach to continuum-limit uniqueness is to
establish Borel summability of the perturbative expansion,
then apply Watson's theorem.  Our argument avoids this:
we follow the RG flow \emph{directly} via the lattice
CS equation rather than reconstructing the theory from
its perturbative expansion.  The lattice provides
deterministic dynamics, and Gronwall's inequality gives
convergence.  Borel summability, if established, would
give an alternative proof but is not required.
\end{remark}

\subsection{Non-triviality of the continuum theory}

We now show the continuum limit is an \emph{interacting}
QFT, not a generalized free field.  A generalized free
field satisfies Wick's theorem: all connected $n$-point
functions vanish for $n \geq 3$.  We prove the connected
four-point function is non-zero.

\begin{definition}[Dimensionless four-point coupling]
For a gauge-invariant local operator $\mathcal{O}$
of dimension~$\Delta_{\mathcal{O}}$
\textup{(}e.g.\ $\mathcal{O} = \tr F_{\mu\nu}^2$\textup{)},
define the renormalized Schwinger functions
$G_R^{(n)} = Z_{\mathcal{O}}^n \cdot G_{\mathrm{bare}}^{(n)}$
and the dimensionless four-point coupling
\begin{equation}\label{eq:lambda4}
  \lambda_4
  \;:=\; \frac{G_{R,\mathrm{conn}}^{(4)}(p_1,\ldots,p_4)}
  {\bigl[G_R^{(2)}(p_1)\bigr]^2}
  \bigg|_{|p_i| = \mu},
\end{equation}
evaluated at a common momentum scale~$\mu$.
\end{definition}

\begin{lemma}[RG invariance of $\lambda_4$]
\label{lem:lambda4_rg}
The ratio $\lambda_4$ is independent of the
renormalization factors~$Z_{\mathcal{O}}$:
since $G_{R,\mathrm{conn}}^{(4)}$ has four operator
insertions and $[G_R^{(2)}]^2$ has $2 \times 2 = 4$
insertions, the $Z_{\mathcal{O}}^4$ factors cancel
in~\eqref{eq:lambda4}.  Consequently, $\lambda_4$
is scheme-independent and satisfies the homogeneous
CS equation
$[a\,\partial_a + \beta_L\,\partial_g]\,\lambda_4 = 0$.
\end{lemma}

\begin{proposition}[Conditional non-triviality]
\label{thm:nontrivial}
Assume (B1)--(B4) \cite{FontesSU2}.
The continuum Euclidean Yang--Mills theory has a
non-vanishing connected four-point function:
$\lambda_4 \neq 0$.  In particular, the theory is not
a generalized free field.
\end{proposition}

\begin{proof}
\textbf{Step~1: Perturbative contribution.}
At a perturbative scale $\mu \gg \Lambda_{\mathrm{QCD}}$,
the connected four-point function of $\tr F^2$ is given
to leading order by one-loop gluon box and crossed
diagrams:
\begin{equation}\label{eq:lambda4_pert}
  \lambda_4^{\mathrm{pert}}(\mu)
  = c_1\,g(\mu)^4 + c_2\,g(\mu)^6 + \cdots
\end{equation}
with $c_1 \neq 0$ (an explicit, computable constant
depending on the external kinematics and the gauge group).
This is a standard perturbative calculation~\cite{Seiler82}.

\smallskip
\textbf{Step~2: Non-perturbative bound.}
By Proposition~\ref{prop:dob_analytic} (Dobrushin
analyticity at the Dobrushin scale), the polymer expansion
bounds~(B1--B4) transfer to individual Schwinger functions.
The non-perturbative contribution satisfies:
\begin{equation}
  |\lambda_4 - \lambda_4^{\mathrm{pert}}|
  \;\leq\; C\,e^{-c/g(\mu)^2}.
\end{equation}
Along the asymptotic freedom trajectory at
scale~$\mu$: $g(\mu)^2 = 1/(2\beta_0 \log(\mu
/\Lambda_{\mathrm{QCD}}))$.  For
$\mu = R\,\Lambda_{\mathrm{QCD}}$ with $R$ large:
\begin{equation}
  e^{-c/g(\mu)^2}
  = e^{-2c\beta_0\,\log R}
  = R^{-2c\beta_0}.
\end{equation}
For $\SU(2)$ with $c = 8\pi^2/N$:
$2c\beta_0 \approx 3.67$, so the non-perturbative
correction decays as $R^{-3.67}$.

\smallskip
\textbf{Step~3: Perturbative dominance.}
Choose $R$ large enough that
$C\,R^{-3.67} < \tfrac{1}{2}\,|c_1|\,g(R\Lambda)^4$.
Since $g(R\Lambda)^4 \sim 1/(\log R)^2$ decays only
logarithmically while $R^{-3.67}$ decays as a power,
such $R$ exists.  At this scale:
\begin{equation}
  |\lambda_4| \;\geq\;
  |c_1|\,g^4 - C\,e^{-c/g^2}
  \;>\; \tfrac{1}{2}\,|c_1|\,g^4 \;>\; 0.
\end{equation}

\smallskip
\textbf{Step~4: Continuum limit.}
By Proposition~\ref{thm:unique}(b), the renormalized
Schwinger functions converge as $a \to 0$.  By
Lemma~\ref{lem:lambda4_rg}, $\lambda_4$ is
RG-invariant (the $Z$-factors cancel).  Therefore the
continuum value equals the lattice value at any scale:
\begin{equation}
  \lambda_4^{\mathrm{cont}}
  = \lim_{a \to 0} \lambda_4^{\mathrm{latt}}(a)
  = \lambda_4(\mu)
  \;\neq\; 0.
\end{equation}
A generalized free field has
$G_{\mathrm{conn}}^{(4)} \equiv 0$, hence
$\lambda_4 = 0$.  Since $\lambda_4 \neq 0$,
the continuum theory is not free.
\end{proof}

\begin{remark}[Alternative argument via the Dobrushin scale]
\label{rem:dobrushin_nontrivial}
The non-triviality can also be seen directly from the
effective theory at the Dobrushin scale~$k_*$.  At this
scale, $\beta_{\mathrm{eff}} \leq \beta_c/L^d \approx 0.002$
for $\SU(2)$ in $d = 4$: the theory is in deep strong
coupling.  The Binder cumulant
$U_4 = -\kappa_4/(3\kappa_2^2)$
(where $\kappa_n$ are cumulants of $\tr F^2$)
satisfies $U_4 \approx 1/3$ at such small $\beta$
(approaching the Haar-measure value),
compared to $U_4 = 0$ for a Gaussian.
The polymer corrections~$R_{k_*}$ are small in the
Dobrushin regime, so $U_4$ remains bounded away from~$0$.
Since $U_4 \propto \lambda_4$, this gives a second proof
of $\lambda_4 \neq 0$.
\end{remark}

\subsection{Reflection positivity and OS reconstruction}

A central ingredient for promoting the Euclidean continuum
theory to a relativistic QFT is \emph{reflection positivity}
(RP).  The following proposition resolves what was originally
identified as hypothesis~(H2).

\begin{proposition}[Reflection positivity of the continuum
limit]\label{prop:rp}
Assume (B1)--(B4) \cite{FontesSU2}.
The continuum Schwinger functions
$G_R^{(n)}(x_1, \ldots, x_n)$ of
Proposition~\ref{thm:unique} satisfy
Osterwalder--Schrader reflection positivity:
for any hyperplane~$\Pi$, any $m \geq 1$, and any
Schwartz test function $f \in \mathcal{S}(\R^{dm})$
supported in the open half-space~$\Pi_+$,
\begin{equation}\label{eq:rp}
  \int f(x_1,\ldots,x_m)\,\overline{f(\theta x_1,
  \ldots,\theta x_m)}\,
  G_R^{(2m)}(\theta x_1, \ldots, \theta x_m,
  x_1, \ldots, x_m)\,
  d^{dm}x \;\geq\; 0,
\end{equation}
where $\theta$ denotes reflection through~$\Pi$.
\end{proposition}

\begin{proof}
The argument has three steps.

\smallskip
\noindent\textbf{Step~1: Lattice RP.}
The Wilson action on the hypercubic lattice satisfies RP
with respect to hyperplanes between adjacent time-slices
(Osterwalder--Seiler~\cite{OsterwalderSeiler78}).
The proof uses the character expansion of the transfer
matrix: $e^{-\beta(1 - (1/N)\Re\tr U)}
= \sum_R c_R(\beta)\,\chi_R(U)$ with $c_R \geq 0$,
which gives a positive kernel across the reflection
plane.  This establishes~\eqref{eq:rp} for the
\emph{bare} lattice Schwinger functions $G_{\mathrm{bare}}^{(2m)}(\cdot\,;a)$
at each lattice spacing~$a$.

\smallskip
\noindent\textbf{Step~2: Positive renormalization.}
The renormalized Schwinger functions are
$G_R^{(n)} = Z_\mathcal{O}(a)^n\,G_{\mathrm{bare}}^{(n)}$
where $Z_\mathcal{O}(a)
= \exp\!\bigl(-n\!\int_a^{a_0}\!\gamma_\mathcal{O}
(g(a'))\,da'/a'\bigr)$.
Since $\gamma_\mathcal{O}$ is real for the Euclidean
theory, $Z_\mathcal{O}(a) \in \R_{>0}$.
The RP integrand in~\eqref{eq:rp} at lattice
spacing~$a$ is therefore
$Z_\mathcal{O}^{2m} \times (\text{bare RP integrand})$
with $Z_\mathcal{O}^{2m} > 0$, so lattice RP
implies~\eqref{eq:rp} for the renormalized Schwinger
functions at each~$a$.

\smallskip
\noindent\textbf{Step~3: Passage to the limit.}
By Proposition~\ref{thm:unique}(b), $G_R^{(2m)}(\cdot\,;a)
\to G_R^{(2m)}(\cdot)$ as $a \to 0$, pointwise at fixed
physical separations.  The uniform bound
\begin{equation}
  |G_R^{(2m)}(\theta x, x;\, a)|
  \;\leq\; C_m\,e^{-\Delta\,\sum_{i<j}|x_i - x_j|}
\end{equation}
holds because: at a reference scale~$a_0$, $G_R(\cdot\,;a_0)$
is finite and exponentially decaying (mass gap); the CS
flow from $a_0$ to $a < a_0$ introduces only $O(a_0^2)$
corrections (Lemma~\ref{lem:irrelevant}), so
$|G_R(\cdot\,;a)| \leq |G_R(\cdot\,;a_0)|\,(1 + C\,a_0^2)$
uniformly for $a \leq a_0$.
\textup{(}Note: the \emph{bare} correlation bound from
Theorem~\ref{thm:conditional} diverges as $a \to 0$
(Remark~\ref{rem:bare_vs_phys}), but the \emph{renormalized}
bound is finite because $Z(a) \to \infty$ compensates
the shrinking bare correlations.\textup{)}
This provides an $L^1$-dominating function
for the integrand $|f|^2 \times |G_R^{(2m)}|$.
By the dominated convergence theorem:
\begin{equation}
  \int |f|^2 \,G_R^{(2m)}(\cdot)
  = \lim_{a \to 0} \int |f|^2\,G_R^{(2m)}(\cdot\,;a)
  \;\geq\; 0,
\end{equation}
since each term in the limit is~$\geq 0$.
Euclidean covariance of the continuum limit (from
universality, Remark~\ref{rem:euclidean}) extends RP
from the coordinate hyperplanes to arbitrary~$\Pi$.
\end{proof}

\begin{remark}[The RG is a proof tool, not a definition]
\label{rem:rg_tool}
A crucial point: the continuum limit is defined as
$\lim_{a \to 0} G_R^{(n)}(\cdot\,;a)$, the limit
of renormalized Schwinger functions of the
\emph{original lattice Wilson theory}.  The block-spin RG
is used only to \emph{prove} that this limit exists
(via the mass gap, Dobrushin uniqueness, and CS
convergence).  The effective measures at intermediate RG
scales are analytical scaffolding; they need not satisfy
RP themselves.  Reflection positivity is inherited
\emph{directly} from the lattice Wilson action to the
continuum limit, bypassing the RG entirely.
\end{remark}

\begin{remark}[Euclidean covariance]
\label{rem:euclidean}
The lattice has hypercubic symmetry, while the continuum
theory should have full $O(d)$ Euclidean covariance.
By Lemma~\ref{lem:irrelevant}, the rotation-breaking
lattice artifacts are $O(a^2)$ (Symanzik improvement);
they vanish in the continuum limit.  The uniqueness of
the limit (Proposition~\ref{thm:unique}(c)) then implies
full $O(d)$ covariance, since the continuum theory cannot
``remember'' the orientation of the lattice.
\end{remark}

\begin{remark}[OS axioms and reconstruction]
\label{rem:os}
With reflection positivity established, the full set of
Osterwalder--Schrader axioms is
verified~\cite{OsterwalderSchrader}:
\textup{(OS0)}~temperedness (from exponential decay);
\textup{(OS1)}~Euclidean covariance (Remark~\ref{rem:euclidean});
\textup{(OS2)}~reflection positivity (Proposition~\ref{prop:rp});
\textup{(OS3)}~symmetry (bosonic fields);
\textup{(OS4)}~cluster property (exponential clustering from
mass gap).
By the OS reconstruction theorem, there exists a
Wightman QFT with positive Hamiltonian, unique vacuum, and
mass gap $\Delta > 0$.  This completes the proof of
Proposition~\ref{thm:unique}(e)--(f).
\end{remark}

\subsection{Summary of the continuum theory}

In $d = 2$, the mass gap and Gibbs uniqueness hold
unconditionally for every compact~$G$ and every
$\beta > 0$ (Theorem~\ref{thm:d2}).  For $d \geq 3$,
the following results hold conditional on
(B1)--(B4) \cite{FontesSU2}:

\begin{center}
\renewcommand{\arraystretch}{1.3}
\begin{tabular}{@{}cll@{}}
\toprule
& \textbf{Result} & \textbf{Source} \\
\midrule
B5 & Mass gap $m \cdot a \geq c/\beta^{1/d}$
  & Theorem~\ref{thm:conditional} \\
& Euclidean continuum limit exists
  & Proposition~\ref{thm:unique}(b) \\
& Euclidean limit is unique
  & Proposition~\ref{thm:unique}(c) \\
& Theory determined by $\Lambda_{\mathrm{QCD}}$
  & Proposition~\ref{thm:unique}(c) \\
& Theory is interacting ($\lambda_4 \neq 0$)
  & Proposition~\ref{thm:nontrivial} \\
& Reflection positivity
  & Proposition~\ref{prop:rp} \\
& Wightman axioms with $\Delta > 0$
  & Remark~\ref{rem:os} \\
\bottomrule
\end{tabular}
\end{center}

\noindent
All results in the table above ($d \geq 3$) are conditional
on (B1)--(B4) \cite{FontesSU2} alone.


%% ====================================================================
\section{Discussion}\label{sec:discussion}
%% ====================================================================

\subsection{Architecture of the proof}

The proof of the mass gap
(Theorem~\ref{thm:conditional}) has a two-step
architecture: the effective action control~\cite{FontesSU2}
governs the RG flow
through all coupling scales, and Dobrushin uniqueness closes
the argument at scale~$k_*$.

\begin{center}
\renewcommand{\arraystretch}{1.2}
\begin{tabular}{@{}lcl@{}}
\toprule
\textbf{Scale} & \textbf{Input} &
\textbf{Role} \\ \midrule
$k = 0, \ldots, k_*-1$
  & (B1)--(B4)~\cite{FontesSU2} & RG flow control \\
$k = k_*$: $\beta_{\mathrm{eff}} \leq \beta_c/L^d$
  & Dobrushin uniqueness & $\rho_\Lambda \geq 0.69$ \\
\bottomrule
\end{tabular}
\end{center}

The mass gap bound $m \cdot a \geq c/\beta^{1/d}$ is
controlled entirely by the strong-coupling Dobrushin gap
divided by the number of RG steps:
$c/\beta^{1/d} = \rho_* / L^{k_*}$.

The semi-classical HS gap
($\geq c_1\beta_{\mathrm{eff}}$, Method~A) and
the entropy decay ($W_k$ shrinking, Method~B) are
\emph{not} used in the proof of
Theorem~\ref{thm:conditional}.  They serve a different
purpose: establishing the unconditional FSC
(Theorem~\ref{thm:fsc}), which yields a
$\beta$-independent mass gap at the cost of triviality
(Theorem~\ref{thm:obstruction}).

\subsection{The intermediate coupling regime}

The effective action control~(B1)--(B4) governs the RG flow
through the
intermediate coupling regime $\beta_{\mathrm{eff}} \approx 1$.
This regime consists of $O(1)$ RG steps
(specifically, $\leq 3$ for $d = 4$, $L = 2$).
For $G = \SU(2)$, $d = 4$, the companion
paper~\cite{FontesSU2} verifies (B1)--(B4) via explicit
Hessian spectrum, parity, and cluster expansion.

The FSC result (Theorem~\ref{thm:fsc}) bypasses this
regime entirely by working at a single scale, achieving
an $O(1)$ block gap via the single-well monotonicity
structure.  However, the resulting $\beta$-independent mass
gap leads to a trivial continuum limit
(Remark~\ref{rem:triviality}).

This reveals a fundamental tension: any argument that
avoids the intermediate coupling regime by using only
single-scale information will necessarily produce a
$\beta$-independent bound, because the $\beta$-dependence
of correlations on a fixed lattice is a \emph{multi-scale}
phenomenon.  The correlation length $\xi(\beta)$ grows
because fluctuations at each scale contribute cumulatively,
and tracking this accumulation requires following the RG
through intermediate coupling.

For general compact simple~$G$ beyond $\SU(2)$,
verifying~(B1)--(B4) remains open.
Three potential approaches:

\begin{enumerate}[nosep]
  \item \textbf{Direct verification.}
    Extend the methods of~\cite{FontesSU2} (Hessian spectrum,
    parity, cluster expansion) to general~$G$.  The main
    additional input needed is the representation-theoretic
    analysis of the character expansion and the Christoffel
    bound for the non-abelian Hessian.

  \item \textbf{Entropy-to-pointwise transfer.}
    The entropy bound gives
    $\KL(\mu_{k_*} \| \nu_{k_*}) \leq C/\beta^{2}$,
    which is $\ell^1$ (average) control.
    A concentration inequality or hypercontractivity argument
    could potentially upgrade this to $\ell^\infty$
    (pointwise) control, at least for the per-link
    oscillation.  The per-plaquette variance bound
    (Proposition~\ref{prop:plaquette}) is a step in this
    direction.
    Establishing the geometric decay
    (Conjecture~\ref{conj:entropy_geom})
    would sharpen the bound to $C/\beta^{2+\alpha}$.

  \item \textbf{Convexity-based RG.}
    On the restricted ensemble $\Omega_{\mathrm{good}}$
    where $|\theta_p| \leq M = O(\sqrt{\log\beta/\beta})$
    for all plaquettes~$p$, the Wilson action is
    strictly convex: $\Hess(S_W) \geq (1 - O(\log\beta/\beta))
    \,B^\top B > 0$.  If the RG could be shown to preserve
    this convexity through the $O(1)$ intermediate steps
    (perhaps with a controlled loss), the Brascamp--Lieb
    inequality would give pointwise bounds on the effective
    action without the polymer expansion.  The entropy
    bound ensures the restricted ensemble has probability
    $\geq 1 - O(1/\beta)$ at all weak-coupling scales.
\end{enumerate}

\subsection{Summary of contributions}

\begin{center}
\renewcommand{\arraystretch}{1.3}
\begin{tabular}{@{}p{5.5cm}cl@{}}
\toprule
\textbf{Component} & \textbf{Status} &
\textbf{Source} \\ \midrule
Lattice mass gap (strong coupling)
  & Done & OS~(1978) \\
UV stability (effective action)
  & Proved ($\SU(2)$) & \cite{FontesSU2} \\
Lattice mass gap (weak coupling)
  & \textbf{This paper} & Thm~\ref{thm:conditional} \\
FSC (positive constant)
  & \textbf{This paper}${}^\dagger$ & Thm~\ref{thm:fsc} \\
Unique Gibbs measure
  & \textbf{This paper}${}^\dagger$ & Thm~\ref{thm:fsc} \\
Euclidean continuum limit exists
  & \textbf{This paper} & Prop~\ref{thm:unique} \\
Euclidean limit is unique
  & \textbf{This paper} & Prop~\ref{thm:unique} \\
Non-trivial (interacting)
  & \textbf{This paper} & Prop~\ref{thm:nontrivial} \\
Reflection positivity
  & \textbf{This paper} & Prop~\ref{prop:rp} \\
Wightman axioms with mass gap
  & \textbf{This paper} & Remark~\ref{rem:os} \\
\bottomrule
\end{tabular}
\end{center}

\noindent
For $G = \SU(2)$, $d = 4$: all results are unconditional
(via~\cite{FontesSU2}).
For general compact simple~$G$: all results hold once
the analogous (B1)--(B4) is verified.

\noindent
{\small ${}^\dagger$FSC and Gibbs uniqueness additionally require
the unique minimum conjecture
(Conjecture~\ref{conj:unique_min}), proved for $G = U(1)$.
All other results, including the continuum limit and
Wightman axioms, are independent of the FSC.}

\medskip
The key contributions are:
(i)~the bridge from effective action control
to the mass gap (Theorem~\ref{thm:conditional}),
via Dobrushin uniqueness at the threshold scale;
(ii)~the FSC
(Theorem~\ref{thm:fsc}, conditional on the unique minimum
conjecture), proving Gibbs uniqueness and positive mass gap
at every coupling;
(iii)~the Dobrushin analyticity transfer
(Proposition~\ref{prop:dob_analytic}), which converts
polymer expansion bounds into Schwinger function bounds;
(iv)~the resolution of the Reisz power-counting input
(Lemma~\ref{lem:irrelevant}), establishing $O(a^2)$
irrelevant corrections to the CS equation;
(v)~the CS convergence argument giving uniqueness
without Borel summability (Proposition~\ref{thm:unique});
(vi)~the non-triviality via perturbative dominance
at high scales (Proposition~\ref{thm:nontrivial});
and (vii)~the reflection positivity of the continuum
limit (Proposition~\ref{prop:rp}), inherited from the
lattice Wilson action and enabling OS reconstruction.


%% ====================================================================
\section*{Acknowledgments}
%% ====================================================================

The author thanks Claude (Anthropic) for sustained
collaboration in developing the mathematical framework.

\begin{thebibliography}{99}

\bibitem{FontesSU2} R.~Fontes,
  ``Mass gap for $\SU(2)$ lattice Yang--Mills:
  block-spin RG with explicit Hessian spectrum,
  parity, and Faddeev--Popov factorization,''
  preprint (2026).

\bibitem{Balaban85} T.~Ba{\l}aban,
  ``Ultraviolet stability in field theory:
  the $\phi^4_3$ model,''
  in \emph{Scaling and Self-Similarity in Physics},
  Birkh\"auser (1985).

\bibitem{Balaban89} T.~Ba{\l}aban,
  ``Large field renormalization, II:
  Localization, exponentiation, and bounds for the
  $\R$ operation,''
  Comm.\ Math.\ Phys.\ \textbf{122} (1989), 355--392.
  \texttt{doi:10.1007/BF01238433}

\bibitem{Guth80} A.~H.~Guth,
  ``Existence proof of a nonconfining phase in
  four-dimensional $U(1)$ lattice gauge theory,''
  Phys.\ Rev.\ D \textbf{21} (1980), 2291--2307.
  \texttt{doi:10.1103/PhysRevD.21.2291}

\bibitem{GopfertMack82} M.~G\"opfert and G.~Mack,
  ``Proof of confinement of static quarks in 3-dimensional
  $U(1)$ lattice gauge theory for all values of the
  coupling constant,''
  Comm.\ Math.\ Phys.\ \textbf{82} (1982), 545--606.
  \texttt{doi:10.1007/BF01961240}

\bibitem{HelfferSjostrand} B.~Helffer and J.~Sj\"ostrand,
  ``Multiple wells in the semi-classical limit, I,''
  Comm.\ Partial Differential Equations \textbf{9}
  (1984), 337--408.
  \texttt{doi:10.1080/03605308408820335}

\bibitem{Helffer02} B.~Helffer,
  \emph{Semi-Classical Analysis, Witten Laplacians,
  and Statistical Mechanics},
  Lecture Notes in Math., vol.~1702, Springer, 2002.
  \texttt{doi:10.1007/978-3-662-04936-4}

\bibitem{DobrushinShlosman} R.~L.~Dobrushin and
  S.~B.~Shlosman,
  ``Completely analytical interactions:
  constructive description,''
  J.~Stat.\ Phys.\ \textbf{46} (1987), 983--1014.
  \texttt{doi:10.1007/BF01011153}

\bibitem{Martinelli1999} F.~Martinelli,
  ``Lectures on Glauber dynamics for discrete spin
  models,''
  in \emph{Lectures on Probability Theory and Statistics},
  Lecture Notes in Math., vol.~1717, Springer (1999),
  93--191.
  \texttt{doi:10.1007/978-3-540-48115-7\_2}

\bibitem{HolleyStroock} R.~Holley and D.~Stroock,
  ``Logarithmic Sobolev inequalities and stochastic
  Ising models,''
  J.~Stat.\ Phys.\ \textbf{46} (1987), 1159--1194.
  \texttt{doi:10.1007/BF01011161}

\bibitem{StroockZegarlinski} D.~Stroock and
  B.~Zegarlinski,
  ``The logarithmic Sobolev inequality for continuous
  spin systems on a lattice,''
  J.~Funct.\ Anal.\ \textbf{104} (1992), 299--326.
  \texttt{doi:10.1016/0022-1236(92)90003-2}

\bibitem{Singer78} I.~M.~Singer,
  ``Some remarks on the Gribov ambiguity,''
  Comm.\ Math.\ Phys.\ \textbf{60} (1978), 7--12.
  \texttt{doi:10.1007/BF01609471}

\bibitem{OsterwalderSchrader} K.~Osterwalder and
  R.~Schrader,
  ``Axioms for Euclidean Green's functions,''
  Comm.\ Math.\ Phys.\ \textbf{31} (1973), 83--112;
  \textbf{42} (1975), 281--305.
  \texttt{doi:10.1007/BF01645738}

\bibitem{Seiler82} E.~Seiler,
  \emph{Gauge Theories as a Problem of Constructive
  Quantum Field Theory and Statistical Mechanics},
  Lecture Notes in Phys., vol.~159, Springer, 1982.
  \texttt{doi:10.1007/3-540-11559-5}

\bibitem{OsterwalderSeiler78} K.~Osterwalder and
  E.~Seiler,
  ``Gauge field theories on a lattice,''
  Ann.\ Physics \textbf{110} (1978), 440--471.
  \texttt{doi:10.1016/0003-4916(78)90039-8}

\bibitem{Reisz1988} T.~Reisz,
  ``A power counting theorem for Feynman integrals
  on the lattice,''
  Comm.\ Math.\ Phys.\ \textbf{116} (1988), 81--126.
  \texttt{doi:10.1007/BF01239027}

\bibitem{Reisz89} T.~Reisz,
  ``Renormalization of Feynman integrals on the lattice,''
  Comm.\ Math.\ Phys.\ \textbf{117} (1989), 79--108.
  \texttt{doi:10.1007/BF01228412}

\bibitem{Rothaus} O.~S.~Rothaus,
  ``Diffusion on compact Riemannian manifolds and
  logarithmic Sobolev inequalities,''
  J.~Funct.\ Anal.\ \textbf{42} (1981), 102--109.
  \texttt{doi:10.1016/0022-1236(81)90049-5}

\bibitem{Luscher_Weisz} M.~L\"uscher and P.~Weisz,
  ``On-shell improved lattice gauge theories,''
  Comm.\ Math.\ Phys.\ \textbf{97} (1985), 59--77.
  \texttt{doi:10.1007/BF01206178}

\bibitem{Akhiezer} N.~I.~Akhiezer,
  \emph{The Classical Moment Problem},
  Oliver \& Boyd, Edinburgh, 1965.

\bibitem{Dobrushin68} R.~L.~Dobrushin,
  ``Description of a random field by means of conditional
  probabilities and conditions for its regularity,''
  Theor.\ Probab.\ Appl.\ \textbf{13} (1968), 197--224.
  \texttt{doi:10.1137/1113026}

\bibitem{Simon_statistical} B.~Simon,
  \emph{The Statistical Mechanics of Lattice Gases},
  vol.~I, Princeton Univ.\ Press, 1993.

\end{thebibliography}

\end{document}
